%% Generated by Sphinx.
\def\sphinxdocclass{report}
\documentclass[a4paper,11pt,english]{sphinxmanual}
\ifdefined\pdfpxdimen
   \let\sphinxpxdimen\pdfpxdimen\else\newdimen\sphinxpxdimen
\fi \sphinxpxdimen=.75bp\relax
\ifdefined\pdfimageresolution
    \pdfimageresolution= \numexpr \dimexpr1in\relax/\sphinxpxdimen\relax
\fi
%% let collapsible pdf bookmarks panel have high depth per default
\PassOptionsToPackage{bookmarksdepth=5}{hyperref}

\PassOptionsToPackage{booktabs}{sphinx}
\PassOptionsToPackage{colorrows}{sphinx}

\PassOptionsToPackage{warn}{textcomp}
\usepackage[utf8]{inputenc}
\ifdefined\DeclareUnicodeCharacter
% support both utf8 and utf8x syntaxes
  \ifdefined\DeclareUnicodeCharacterAsOptional
    \def\sphinxDUC#1{\DeclareUnicodeCharacter{"#1}}
  \else
    \let\sphinxDUC\DeclareUnicodeCharacter
  \fi
  \sphinxDUC{00A0}{\nobreakspace}
  \sphinxDUC{2500}{\sphinxunichar{2500}}
  \sphinxDUC{2502}{\sphinxunichar{2502}}
  \sphinxDUC{2514}{\sphinxunichar{2514}}
  \sphinxDUC{251C}{\sphinxunichar{251C}}
  \sphinxDUC{2572}{\textbackslash}
\fi
\usepackage{cmap}
\usepackage[T1]{fontenc}
\usepackage{amsmath,amssymb,amstext}
\usepackage{babel}



\usepackage{tgtermes}
\usepackage{tgheros}
\renewcommand{\ttdefault}{txtt}



\usepackage[Bjarne]{fncychap}
\usepackage[,numfigreset=1,mathnumfig,mathnumsep={.}]{sphinx}
\sphinxsetup{VerbatimColor={rgb}{0.95,0.95,0.95}}
\fvset{fontsize=auto}
\usepackage{geometry}


% Include hyperref last.
\usepackage{hyperref}
% Fix anchor placement for figures with captions.
\usepackage{hypcap}% it must be loaded after hyperref.
% Set up styles of URL: it should be placed after hyperref.
\urlstyle{same}


\usepackage{sphinxmessages}
\setcounter{tocdepth}{2}



\title{GIdisk2obs to RADMC3D}
\date{May 13, 2025}
\release{05\sphinxhyphen{}2025}
\author{Yi\sphinxhyphen{}Ping Hung}
\newcommand{\sphinxlogo}{\vbox{}}
\renewcommand{\releasename}{Release}
\makeindex
\begin{document}

\ifdefined\shorthandoff
  \ifnum\catcode`\=\string=\active\shorthandoff{=}\fi
  \ifnum\catcode`\"=\active\shorthandoff{"}\fi
\fi

\pagestyle{empty}
\sphinxmaketitle
\pagestyle{plain}
\sphinxtableofcontents
\pagestyle{normal}
\phantomsection\label{\detokenize{index::doc}}


\sphinxstepscope


\chapter{Introduction}
\label{\detokenize{introduction:introduction}}\label{\detokenize{introduction::doc}}

\section{What is this?}
\label{\detokenize{introduction:what-is-this}}
\sphinxAtStartPar
“GIdisk2obs to RADMC3D Pipeline” is a Python package that provides a pipeline for
converting GIdisk2obs data to RADMC3D format. It is designed to facilitate the process of
converting GIdisk2obs data into a format that can be used with RADMC3D, a popular
radiative transfer code used in astrophysics. The package includes a set of tools and
functions that allow users to easily convert GIdisk2obs data into RADMC3D format, as well
as to visualize and analyze the data. The package is designed to be user\sphinxhyphen{}friendly and
easy to use, making it accessible to researchers and scientists working in the field of
astrophysics. The package is built on top of the GIdisk2obs and RADMC3D software
packages, and is designed to work seamlessly with these tools. The package is open\sphinxhyphen{}source
and is available on GitHub, allowing users to contribute to its development and
improvement. The package is actively maintained and updated, with new features and
improvements being added regularly. Overall, “GIdisk2obs to RADMC3D Pipeline” is a
powerful and flexible tool for researchers and scientists working in the field of
astrophysics, providing a simple and efficient way to convert GIdisk2obs data into
RADMC3D format and to visualize and analyze the data. The package is designed to be
easy to use and accessible to researchers and scientists working in the field of
astrophysics, making it a valuable tool for anyone working with GIdisk2obs and RADMC3D
data.


\section{Capabilities}
\label{\detokenize{introduction:capabilities}}

\section{Requirements}
\label{\detokenize{introduction:requirements}}\begin{itemize}
\item {} 
\sphinxAtStartPar
Python

\item {} 
\sphinxAtStartPar
RADMC\sphinxhyphen{}3D

\item {} 
\sphinxAtStartPar
NumPy

\item {} 
\sphinxAtStartPar
Matplotlib

\item {} 
\sphinxAtStartPar
SciPy

\item {} 
\sphinxAtStartPar
Astropy

\end{itemize}

\sphinxstepscope


\chapter{X22 Model}
\label{\detokenize{x22:x22-model}}\label{\detokenize{x22::doc}}

\section{Model Assumptions}
\label{\detokenize{x22:model-assumptions}}

\section{Modules}
\label{\detokenize{x22:modules}}\index{generate\_opacity\_table\_x22() (in module disk\_model)@\spxentry{generate\_opacity\_table\_x22()}\spxextra{in module disk\_model}}

\begin{fulllineitems}
\phantomsection\label{\detokenize{x22:disk_model.generate_opacity_table_x22}}
\pysigstartsignatures
\pysiglinewithargsret
{\sphinxcode{\sphinxupquote{disk\_model.}}\sphinxbfcode{\sphinxupquote{generate\_opacity\_table\_x22}}}
{\sphinxparam{\DUrole{n}{a\_min}}\sphinxparamcomma \sphinxparam{\DUrole{n}{a\_max}}\sphinxparamcomma \sphinxparam{\DUrole{n}{q}}\sphinxparamcomma \sphinxparam{\DUrole{n}{dust\_to\_gas}}\sphinxparamcomma \sphinxparam{\DUrole{n}{precomputed\_grain\_properties\_fname}\DUrole{o}{=}\DUrole{default_value}{\textquotesingle{}precomputed\_grain\_properties/grain\_properties.pkl\textquotesingle{}}}\sphinxparamcomma \sphinxparam{\DUrole{n}{T\_min}\DUrole{o}{=}\DUrole{default_value}{20}}\sphinxparamcomma \sphinxparam{\DUrole{n}{T\_max}\DUrole{o}{=}\DUrole{default_value}{2000}}\sphinxparamcomma \sphinxparam{\DUrole{n}{N\_T}\DUrole{o}{=}\DUrole{default_value}{50}}}
{}
\pysigstopsignatures
\sphinxAtStartPar
Generate an opacity table for given grain size distribution.
\begin{quote}\begin{description}
\sphinxlineitem{Parameters}\begin{itemize}
\item {} 
\sphinxAtStartPar
\sphinxstyleliteralstrong{\sphinxupquote{a\_min}} (\sphinxstyleliteralemphasis{\sphinxupquote{float}}) \textendash{} Min/Max grain size in grain size distribution

\item {} 
\sphinxAtStartPar
\sphinxstyleliteralstrong{\sphinxupquote{a\_max}} (\sphinxstyleliteralemphasis{\sphinxupquote{float}}) \textendash{} Min/Max grain size in grain size distribution

\item {} 
\sphinxAtStartPar
\sphinxstyleliteralstrong{\sphinxupquote{q}} (\sphinxstyleliteralemphasis{\sphinxupquote{float}}) \textendash{} Slope for grain size distribution

\item {} 
\sphinxAtStartPar
\sphinxstyleliteralstrong{\sphinxupquote{dust\_to\_gas}} (\sphinxstyleliteralemphasis{\sphinxupquote{float}}) \textendash{} dust\sphinxhyphen{}to\sphinxhyphen{}gas mass ratio (before sublimation)

\item {} 
\sphinxAtStartPar
\sphinxstyleliteralstrong{\sphinxupquote{precomputed\_grain\_properties\_fname}} (\sphinxstyleliteralemphasis{\sphinxupquote{str}}) \textendash{} Name of the pkl file storing pre\sphinxhyphen{}computed grain material properties

\item {} 
\sphinxAtStartPar
\sphinxstyleliteralstrong{\sphinxupquote{T\_min}} (\sphinxstyleliteralemphasis{\sphinxupquote{float}}) \textendash{} Min/Max temprature for temprature grid

\item {} 
\sphinxAtStartPar
\sphinxstyleliteralstrong{\sphinxupquote{T\_max}} (\sphinxstyleliteralemphasis{\sphinxupquote{float}}) \textendash{} Min/Max temprature for temprature grid

\item {} 
\sphinxAtStartPar
\sphinxstyleliteralstrong{\sphinxupquote{N\_T}} (\sphinxstyleliteralemphasis{\sphinxupquote{int}}) \textendash{} Temperature grid resolution

\end{itemize}

\sphinxlineitem{Returns}
\sphinxAtStartPar
\begin{itemize}
\item {} 
\sphinxAtStartPar
\sphinxstylestrong{T\_crit} (\sphinxstyleemphasis{numpy array}) \textendash{} Sublimation temperatures

\item {} 
\sphinxAtStartPar
\sphinxstylestrong{T} (\sphinxstyleemphasis{numpy array}) \textendash{} Temperature grid

\item {} 
\sphinxAtStartPar
\sphinxstylestrong{lam} (\sphinxstyleemphasis{numpy array}) \textendash{} Wavelength grid

\item {} 
\sphinxAtStartPar
\sphinxstylestrong{kappa} (\sphinxstyleemphasis{numpy array}) \textendash{} absorption opacity at given lam
2D array, first dimension corresponds to temperature range

\item {} 
\sphinxAtStartPar
\sphinxstylestrong{kappa\_p} (\sphinxstyleemphasis{numpy array}) \textendash{} Planck mean opacity

\item {} 
\sphinxAtStartPar
\sphinxstylestrong{kappa\_r} (\sphinxstyleemphasis{numpy array}) \textendash{} Rosseland mean opacity

\item {} 
\sphinxAtStartPar
\sphinxstylestrong{kappa\_s, kappa\_s\_p, kappa\_s\_r} (\sphinxstyleemphasis{numpy array}) \textendash{} same as kappa, kappa\_p, kappa\_r, but for effective scattering opacity
(1\sphinxhyphen{}g)*kappa\_scatter

\end{itemize}


\end{description}\end{quote}

\end{fulllineitems}

\index{generate\_opacity\_table\_opt() (in module disk\_model)@\spxentry{generate\_opacity\_table\_opt()}\spxextra{in module disk\_model}}

\begin{fulllineitems}
\phantomsection\label{\detokenize{x22:disk_model.generate_opacity_table_opt}}
\pysigstartsignatures
\pysiglinewithargsret
{\sphinxcode{\sphinxupquote{disk\_model.}}\sphinxbfcode{\sphinxupquote{generate\_opacity\_table\_opt}}}
{\sphinxparam{\DUrole{n}{a\_min}}\sphinxparamcomma \sphinxparam{\DUrole{n}{a\_max}}\sphinxparamcomma \sphinxparam{\DUrole{n}{q}}\sphinxparamcomma \sphinxparam{\DUrole{n}{dust\_to\_gas}}\sphinxparamcomma \sphinxparam{\DUrole{n}{T\_min}\DUrole{o}{=}\DUrole{default_value}{20}}\sphinxparamcomma \sphinxparam{\DUrole{n}{T\_max}\DUrole{o}{=}\DUrole{default_value}{2000}}\sphinxparamcomma \sphinxparam{\DUrole{n}{N\_T}\DUrole{o}{=}\DUrole{default_value}{100}}\sphinxparamcomma \sphinxparam{\DUrole{n}{save\_table}\DUrole{o}{=}\DUrole{default_value}{True}}}
{}
\pysigstopsignatures
\sphinxAtStartPar
Generate an opacity table for given grain size distribution.
\begin{quote}\begin{description}
\sphinxlineitem{Parameters}\begin{itemize}
\item {} 
\sphinxAtStartPar
\sphinxstyleliteralstrong{\sphinxupquote{a\_min}} (\sphinxstyleliteralemphasis{\sphinxupquote{float}}) \textendash{} Min/Max grain size in grain size distribution

\item {} 
\sphinxAtStartPar
\sphinxstyleliteralstrong{\sphinxupquote{a\_max}} (\sphinxstyleliteralemphasis{\sphinxupquote{float}}) \textendash{} Min/Max grain size in grain size distribution

\item {} 
\sphinxAtStartPar
\sphinxstyleliteralstrong{\sphinxupquote{q}} (\sphinxstyleliteralemphasis{\sphinxupquote{float}}) \textendash{} Slope for grain size distribution

\item {} 
\sphinxAtStartPar
\sphinxstyleliteralstrong{\sphinxupquote{dust\_to\_gas}} (\sphinxstyleliteralemphasis{\sphinxupquote{float}}) \textendash{} dust\sphinxhyphen{}to\sphinxhyphen{}gas mass ratio (before sublimation)

\item {} 
\sphinxAtStartPar
\sphinxstyleliteralstrong{\sphinxupquote{precomputed\_grain\_properties\_fname}} (\sphinxstyleliteralemphasis{\sphinxupquote{str}}) \textendash{} Name of the pkl file storing pre\sphinxhyphen{}computed grain material properties

\item {} 
\sphinxAtStartPar
\sphinxstyleliteralstrong{\sphinxupquote{T\_min}} (\sphinxstyleliteralemphasis{\sphinxupquote{float}}) \textendash{} Min/Max temprature for temprature grid

\item {} 
\sphinxAtStartPar
\sphinxstyleliteralstrong{\sphinxupquote{T\_max}} (\sphinxstyleliteralemphasis{\sphinxupquote{float}}) \textendash{} Min/Max temprature for temprature grid

\item {} 
\sphinxAtStartPar
\sphinxstyleliteralstrong{\sphinxupquote{N\_T}} (\sphinxstyleliteralemphasis{\sphinxupquote{int}}) \textendash{} Temperature grid resolution

\end{itemize}

\sphinxlineitem{Returns}
\sphinxAtStartPar
\begin{itemize}
\item {} 
\sphinxAtStartPar
\sphinxstylestrong{T\_crit} (\sphinxstyleemphasis{numpy array}) \textendash{} Sublimation temperatures

\item {} 
\sphinxAtStartPar
\sphinxstylestrong{T} (\sphinxstyleemphasis{numpy array}) \textendash{} Temperature grid

\item {} 
\sphinxAtStartPar
\sphinxstylestrong{lam} (\sphinxstyleemphasis{numpy array}) \textendash{} Wavelength grid

\item {} 
\sphinxAtStartPar
\sphinxstylestrong{kappa} (\sphinxstyleemphasis{numpy array}) \textendash{} absorption opacity at given lam
2D array, first dimension corresponds to temperature range

\item {} 
\sphinxAtStartPar
\sphinxstylestrong{kappa\_p} (\sphinxstyleemphasis{numpy array}) \textendash{} Planck mean opacity

\item {} 
\sphinxAtStartPar
\sphinxstylestrong{kappa\_r} (\sphinxstyleemphasis{numpy array}) \textendash{} Rosseland mean opacity

\item {} 
\sphinxAtStartPar
\sphinxstylestrong{kappa\_s, kappa\_s\_p, kappa\_s\_r} (\sphinxstyleemphasis{numpy array}) \textendash{} same as kappa, kappa\_p, kappa\_r, but for effective scattering opacity
(1\sphinxhyphen{}g)*kappa\_scatter

\end{itemize}


\end{description}\end{quote}

\end{fulllineitems}

\index{generate\_disk\_property\_table() (in module disk\_model)@\spxentry{generate\_disk\_property\_table()}\spxextra{in module disk\_model}}

\begin{fulllineitems}
\phantomsection\label{\detokenize{x22:disk_model.generate_disk_property_table}}
\pysigstartsignatures
\pysiglinewithargsret
{\sphinxcode{\sphinxupquote{disk\_model.}}\sphinxbfcode{\sphinxupquote{generate\_disk\_property\_table}}}
{\sphinxparam{\DUrole{n}{opacity\_table}}\sphinxparamcomma \sphinxparam{\DUrole{n}{N\_T\_eff}\DUrole{o}{=}\DUrole{default_value}{30}}\sphinxparamcomma \sphinxparam{\DUrole{n}{N\_Sigma\_cs}\DUrole{o}{=}\DUrole{default_value}{30}}\sphinxparamcomma \sphinxparam{\DUrole{n}{T\_eff\_min}\DUrole{o}{=}\DUrole{default_value}{0.5}}}
{}
\pysigstopsignatures
\sphinxAtStartPar
Generate a table that maps T\_eff and Sigma/cs to other local disk
properties.
\begin{quote}\begin{description}
\sphinxlineitem{Parameters}\begin{itemize}
\item {} 
\sphinxAtStartPar
\sphinxstyleliteralstrong{\sphinxupquote{opacity\_table}} (\sphinxstyleliteralemphasis{\sphinxupquote{dict}}) \textendash{} Opacity table from generate\_opacity\_table()

\item {} 
\sphinxAtStartPar
\sphinxstyleliteralstrong{\sphinxupquote{N\_T\_eff}} (\sphinxstyleliteralemphasis{\sphinxupquote{int}}) \textendash{} Resolution of the grid

\item {} 
\sphinxAtStartPar
\sphinxstyleliteralstrong{\sphinxupquote{N\_Sigma\_cs}} (\sphinxstyleliteralemphasis{\sphinxupquote{int}}) \textendash{} Resolution of the grid

\item {} 
\sphinxAtStartPar
\sphinxstyleliteralstrong{\sphinxupquote{T\_eff\_min}} (\sphinxstyleliteralemphasis{\sphinxupquote{float}}) \textendash{} Minimum T\_eff (max T\_eff is determined automatically)

\end{itemize}

\sphinxlineitem{Returns}
\sphinxAtStartPar
\begin{itemize}
\item {} 
\sphinxAtStartPar
\sphinxstyleemphasis{a dictionary containing}

\item {} 
\sphinxAtStartPar
\sphinxstylestrong{T\_eff\_grid} (\sphinxstyleemphasis{numpy array}) \textendash{} Effective temperature grid

\item {} 
\sphinxAtStartPar
\sphinxstylestrong{Sigma\_cs\_l, Sigma\_cs\_r} (\sphinxstyleemphasis{numpy array}) \textendash{} Min/Max Sigma/cs allowed at given T\_eff

\item {} 
\sphinxAtStartPar
\sphinxstylestrong{x} (\sphinxstyleemphasis{numpy array}) \textendash{} Normalized log(Sigma/cs) grid.
We map x={[}0,1{]} to the full range of allowed log(Sigma/cs)
uniformly.

\item {} 
\sphinxAtStartPar
\sphinxstylestrong{Sigma, cs, tau\_p\_mid, tau\_r\_mid, T\_mid} (\sphinxstyleemphasis{numpy array}) \textendash{} local disk properties
for given (T,x)

\end{itemize}


\end{description}\end{quote}

\end{fulllineitems}

\index{DiskModel (class in disk\_model)@\spxentry{DiskModel}\spxextra{class in disk\_model}}

\begin{fulllineitems}
\phantomsection\label{\detokenize{x22:disk_model.DiskModel}}
\pysigstartsignatures
\pysiglinewithargsret
{\sphinxbfcode{\sphinxupquote{\DUrole{k}{class}\DUrole{w}{ }}}\sphinxcode{\sphinxupquote{disk\_model.}}\sphinxbfcode{\sphinxupquote{DiskModel}}}
{\sphinxparam{\DUrole{n}{opacity\_table}}\sphinxparamcomma \sphinxparam{\DUrole{n}{disk\_property\_table}}}
{}
\pysigstopsignatures
\sphinxAtStartPar
Bases: \sphinxcode{\sphinxupquote{object}}

\sphinxAtStartPar
Parametrized disk model for generating radial porfiles of disk
properties and flux density at given wavelengths.
\index{M (disk\_model.DiskModel attribute)@\spxentry{M}\spxextra{disk\_model.DiskModel attribute}}

\begin{fulllineitems}
\phantomsection\label{\detokenize{x22:disk_model.DiskModel.M}}
\pysigstartsignatures
\pysigline
{\sphinxbfcode{\sphinxupquote{M}}}
\pysigstopsignatures
\sphinxAtStartPar
Total mass (g)
\begin{quote}\begin{description}
\sphinxlineitem{Type}
\sphinxAtStartPar
float

\end{description}\end{quote}

\end{fulllineitems}

\index{Mstar (disk\_model.DiskModel attribute)@\spxentry{Mstar}\spxextra{disk\_model.DiskModel attribute}}

\begin{fulllineitems}
\phantomsection\label{\detokenize{x22:disk_model.DiskModel.Mstar}}
\pysigstartsignatures
\pysigline
{\sphinxbfcode{\sphinxupquote{Mstar}}}
\pysigstopsignatures
\sphinxAtStartPar
Stellar mass (g)
\begin{quote}\begin{description}
\sphinxlineitem{Type}
\sphinxAtStartPar
float

\end{description}\end{quote}

\end{fulllineitems}

\index{Mdot (disk\_model.DiskModel attribute)@\spxentry{Mdot}\spxextra{disk\_model.DiskModel attribute}}

\begin{fulllineitems}
\phantomsection\label{\detokenize{x22:disk_model.DiskModel.Mdot}}
\pysigstartsignatures
\pysigline
{\sphinxbfcode{\sphinxupquote{Mdot}}}
\pysigstopsignatures
\sphinxAtStartPar
Accretion rate (g/s)
\begin{quote}\begin{description}
\sphinxlineitem{Type}
\sphinxAtStartPar
float

\end{description}\end{quote}

\end{fulllineitems}

\index{Rd (disk\_model.DiskModel attribute)@\spxentry{Rd}\spxextra{disk\_model.DiskModel attribute}}

\begin{fulllineitems}
\phantomsection\label{\detokenize{x22:disk_model.DiskModel.Rd}}
\pysigstartsignatures
\pysigline
{\sphinxbfcode{\sphinxupquote{Rd}}}
\pysigstopsignatures
\sphinxAtStartPar
Disk size (cm)
\begin{quote}\begin{description}
\sphinxlineitem{Type}
\sphinxAtStartPar
float

\end{description}\end{quote}

\end{fulllineitems}

\index{Q (disk\_model.DiskModel attribute)@\spxentry{Q}\spxextra{disk\_model.DiskModel attribute}}

\begin{fulllineitems}
\phantomsection\label{\detokenize{x22:disk_model.DiskModel.Q}}
\pysigstartsignatures
\pysigline
{\sphinxbfcode{\sphinxupquote{Q}}}
\pysigstopsignatures
\sphinxAtStartPar
Toomre Q
\begin{quote}\begin{description}
\sphinxlineitem{Type}
\sphinxAtStartPar
float

\end{description}\end{quote}

\end{fulllineitems}

\index{R (disk\_model.DiskModel attribute)@\spxentry{R}\spxextra{disk\_model.DiskModel attribute}}

\begin{fulllineitems}
\phantomsection\label{\detokenize{x22:disk_model.DiskModel.R}}
\pysigstartsignatures
\pysigline
{\sphinxbfcode{\sphinxupquote{R}}}
\pysigstopsignatures
\sphinxAtStartPar
Radius grid (R{[}0{]}=0) (cm)
\begin{quote}\begin{description}
\sphinxlineitem{Type}
\sphinxAtStartPar
numpy array

\end{description}\end{quote}

\end{fulllineitems}



\begin{fulllineitems}

\pysigstartsignatures
\pysigline
{\sphinxbfcode{\sphinxupquote{Sigma,~T\_mid,~tau\_p\_mid,~tau\_r\_mid}}}
\pysigstopsignatures
\sphinxAtStartPar
Radial profile at R{[}1:{]}
\begin{quote}\begin{description}
\sphinxlineitem{Type}
\sphinxAtStartPar
numpy array

\end{description}\end{quote}

\end{fulllineitems}

\index{MR (disk\_model.DiskModel attribute)@\spxentry{MR}\spxextra{disk\_model.DiskModel attribute}}

\begin{fulllineitems}
\phantomsection\label{\detokenize{x22:disk_model.DiskModel.MR}}
\pysigstartsignatures
\pysigline
{\sphinxbfcode{\sphinxupquote{MR}}}
\pysigstopsignatures
\sphinxAtStartPar
M(\textless{}R) profile at R{[}1:{]} (cm)
\begin{quote}\begin{description}
\sphinxlineitem{Type}
\sphinxAtStartPar
numpy array

\end{description}\end{quote}

\end{fulllineitems}

\subsubsection*{Methods}


\begin{savenotes}
\sphinxatlongtablestart
\sphinxthistablewithglobalstyle
\sphinxthistablewithnovlinesstyle
\makeatletter
  \LTleft \@totalleftmargin plus1fill
  \LTright\dimexpr\columnwidth-\@totalleftmargin-\linewidth\relax plus1fill
\makeatother
\begin{longtable}{\X{1}{2}\X{1}{2}}
\sphinxtoprule
\endfirsthead

\multicolumn{2}{c}{\sphinxnorowcolor
    \makebox[0pt]{\sphinxtablecontinued{\tablename\ \thetable{} \textendash{} continued from previous page}}%
}\\
\sphinxtoprule
\endhead

\sphinxbottomrule
\multicolumn{2}{r}{\sphinxnorowcolor
    \makebox[0pt][r]{\sphinxtablecontinued{continues on next page}}%
}\\
\endfoot

\endlastfoot
\sphinxtableatstartofbodyhook

\sphinxAtStartPar
{\hyperref[\detokenize{x22:disk_model.DiskModel.generate_disk_profile}]{\sphinxcrossref{\sphinxcode{\sphinxupquote{generate\_disk\_profile}}}}}({[}Mstar, Mdot, Rd, Q, ...{]})
&
\sphinxAtStartPar
Generate radial disk profile (Sigma, T\_mid, tau\_p\_mid, tau\_r\_mid)
\\
\sphinxhline
\sphinxAtStartPar
{\hyperref[\detokenize{x22:disk_model.DiskModel.generate_observed_flux}]{\sphinxcrossref{\sphinxcode{\sphinxupquote{generate\_observed\_flux}}}}}(cosI{[}, scattering{]})
&
\sphinxAtStartPar
Generate flux density at observed wavelengths
\\
\sphinxhline
\sphinxAtStartPar
{\hyperref[\detokenize{x22:disk_model.DiskModel.set_lam_obs_list}]{\sphinxcrossref{\sphinxcode{\sphinxupquote{set\_lam\_obs\_list}}}}}(lam\_obs\_list)
&
\sphinxAtStartPar
Set wavelengths of observation
\\
\sphinxbottomrule
\end{longtable}
\sphinxtableafterendhook
\sphinxatlongtableend
\end{savenotes}


\begin{savenotes}\sphinxattablestart
\sphinxthistablewithglobalstyle
\centering
\begin{tabulary}{\linewidth}[t]{TT}
\sphinxtoprule
\sphinxtableatstartofbodyhook
\sphinxAtStartPar
\sphinxstylestrong{generate\_observed\_flux\_single\_wavelength}
&\\
\sphinxhline
\sphinxAtStartPar
\sphinxstylestrong{generate\_observed\_flux\_single\_wavelength\_no\_scattering}
&\\
\sphinxhline
\sphinxAtStartPar
\sphinxstylestrong{load\_disk\_property\_functions}
&\\
\sphinxhline
\sphinxAtStartPar
\sphinxstylestrong{load\_opacity\_functions}
&\\
\sphinxhline
\sphinxAtStartPar
\sphinxstylestrong{solve\_local\_disk\_properties}
&\\
\sphinxbottomrule
\end{tabulary}
\sphinxtableafterendhook\par
\sphinxattableend\end{savenotes}
\index{generate\_disk\_profile() (disk\_model.DiskModel method)@\spxentry{generate\_disk\_profile()}\spxextra{disk\_model.DiskModel method}}

\begin{fulllineitems}
\phantomsection\label{\detokenize{x22:disk_model.DiskModel.generate_disk_profile}}
\pysigstartsignatures
\pysiglinewithargsret
{\sphinxbfcode{\sphinxupquote{generate\_disk\_profile}}}
{\sphinxparam{\DUrole{n}{Mstar}\DUrole{o}{=}\DUrole{default_value}{None}}\sphinxparamcomma \sphinxparam{\DUrole{n}{Mdot}\DUrole{o}{=}\DUrole{default_value}{None}}\sphinxparamcomma \sphinxparam{\DUrole{n}{Rd}\DUrole{o}{=}\DUrole{default_value}{None}}\sphinxparamcomma \sphinxparam{\DUrole{n}{Q}\DUrole{o}{=}\DUrole{default_value}{None}}\sphinxparamcomma \sphinxparam{\DUrole{n}{N\_R}\DUrole{o}{=}\DUrole{default_value}{50}}\sphinxparamcomma \sphinxparam{\DUrole{n}{N\_itr}\DUrole{o}{=}\DUrole{default_value}{8}}\sphinxparamcomma \sphinxparam{\DUrole{n}{plot\_itr}\DUrole{o}{=}\DUrole{default_value}{False}}}
{}
\pysigstopsignatures
\sphinxAtStartPar
Generate radial disk profile (Sigma, T\_mid, tau\_p\_mid, tau\_r\_mid)
\begin{quote}\begin{description}
\sphinxlineitem{Parameters}\begin{itemize}
\item {} 
\sphinxAtStartPar
\sphinxstyleliteralstrong{\sphinxupquote{Mstar}} (\sphinxstyleliteralemphasis{\sphinxupquote{float}}) \textendash{} Stellar mass (g)

\item {} 
\sphinxAtStartPar
\sphinxstyleliteralstrong{\sphinxupquote{Mdot}} (\sphinxstyleliteralemphasis{\sphinxupquote{float}}) \textendash{} Accretion rate (g/s)

\item {} 
\sphinxAtStartPar
\sphinxstyleliteralstrong{\sphinxupquote{Rd}} (\sphinxstyleliteralemphasis{\sphinxupquote{float}}) \textendash{} Disk size (cm)

\item {} 
\sphinxAtStartPar
\sphinxstyleliteralstrong{\sphinxupquote{Q}} (\sphinxstyleliteralemphasis{\sphinxupquote{float}}) \textendash{} Toomre Q

\item {} 
\sphinxAtStartPar
\sphinxstyleliteralstrong{\sphinxupquote{N\_R}} (\sphinxstyleliteralemphasis{\sphinxupquote{int}}) \textendash{} Number of R grid points

\end{itemize}

\end{description}\end{quote}

\end{fulllineitems}

\index{generate\_observed\_flux() (disk\_model.DiskModel method)@\spxentry{generate\_observed\_flux()}\spxextra{disk\_model.DiskModel method}}

\begin{fulllineitems}
\phantomsection\label{\detokenize{x22:disk_model.DiskModel.generate_observed_flux}}
\pysigstartsignatures
\pysiglinewithargsret
{\sphinxbfcode{\sphinxupquote{generate\_observed\_flux}}}
{\sphinxparam{\DUrole{n}{cosI}}\sphinxparamcomma \sphinxparam{\DUrole{n}{scattering}\DUrole{o}{=}\DUrole{default_value}{True}}\sphinxparamcomma \sphinxparam{\DUrole{o}{**}\DUrole{n}{kwargs}}}
{}
\pysigstopsignatures
\sphinxAtStartPar
Generate flux density at observed wavelengths

\end{fulllineitems}

\index{generate\_observed\_flux\_single\_wavelength() (disk\_model.DiskModel method)@\spxentry{generate\_observed\_flux\_single\_wavelength()}\spxextra{disk\_model.DiskModel method}}

\begin{fulllineitems}
\phantomsection\label{\detokenize{x22:disk_model.DiskModel.generate_observed_flux_single_wavelength}}
\pysigstartsignatures
\pysiglinewithargsret
{\sphinxbfcode{\sphinxupquote{generate\_observed\_flux\_single\_wavelength}}}
{\sphinxparam{\DUrole{n}{cosI}}\sphinxparamcomma \sphinxparam{\DUrole{n}{lam\_obs}}\sphinxparamcomma \sphinxparam{\DUrole{n}{f\_kappa\_obs}}\sphinxparamcomma \sphinxparam{\DUrole{n}{f\_kappa\_s\_obs}}\sphinxparamcomma \sphinxparam{\DUrole{n}{dtau}\DUrole{o}{=}\DUrole{default_value}{0.2}}}
{}
\pysigstopsignatures
\end{fulllineitems}

\index{generate\_observed\_flux\_single\_wavelength\_no\_scattering() (disk\_model.DiskModel method)@\spxentry{generate\_observed\_flux\_single\_wavelength\_no\_scattering()}\spxextra{disk\_model.DiskModel method}}

\begin{fulllineitems}
\phantomsection\label{\detokenize{x22:disk_model.DiskModel.generate_observed_flux_single_wavelength_no_scattering}}
\pysigstartsignatures
\pysiglinewithargsret
{\sphinxbfcode{\sphinxupquote{generate\_observed\_flux\_single\_wavelength\_no\_scattering}}}
{\sphinxparam{\DUrole{n}{cosI}}\sphinxparamcomma \sphinxparam{\DUrole{n}{lam\_obs}}\sphinxparamcomma \sphinxparam{\DUrole{n}{f\_kappa\_obs}}\sphinxparamcomma \sphinxparam{\DUrole{n}{N\_tau}\DUrole{o}{=}\DUrole{default_value}{50}}\sphinxparamcomma \sphinxparam{\DUrole{n}{tau\_min}\DUrole{o}{=}\DUrole{default_value}{0.01}}}
{}
\pysigstopsignatures
\end{fulllineitems}

\index{load\_disk\_property\_functions() (disk\_model.DiskModel method)@\spxentry{load\_disk\_property\_functions()}\spxextra{disk\_model.DiskModel method}}

\begin{fulllineitems}
\phantomsection\label{\detokenize{x22:disk_model.DiskModel.load_disk_property_functions}}
\pysigstartsignatures
\pysiglinewithargsret
{\sphinxbfcode{\sphinxupquote{load\_disk\_property\_functions}}}
{\sphinxparam{\DUrole{n}{disk\_property\_table}}}
{}
\pysigstopsignatures
\end{fulllineitems}

\index{load\_opacity\_functions() (disk\_model.DiskModel method)@\spxentry{load\_opacity\_functions()}\spxextra{disk\_model.DiskModel method}}

\begin{fulllineitems}
\phantomsection\label{\detokenize{x22:disk_model.DiskModel.load_opacity_functions}}
\pysigstartsignatures
\pysiglinewithargsret
{\sphinxbfcode{\sphinxupquote{load\_opacity\_functions}}}
{\sphinxparam{\DUrole{n}{opacity\_table}}}
{}
\pysigstopsignatures
\end{fulllineitems}

\index{set\_lam\_obs\_list() (disk\_model.DiskModel method)@\spxentry{set\_lam\_obs\_list()}\spxextra{disk\_model.DiskModel method}}

\begin{fulllineitems}
\phantomsection\label{\detokenize{x22:disk_model.DiskModel.set_lam_obs_list}}
\pysigstartsignatures
\pysiglinewithargsret
{\sphinxbfcode{\sphinxupquote{set\_lam\_obs\_list}}}
{\sphinxparam{\DUrole{n}{lam\_obs\_list}}}
{}
\pysigstopsignatures
\sphinxAtStartPar
Set wavelengths of observation

\end{fulllineitems}

\index{solve\_local\_disk\_properties() (disk\_model.DiskModel method)@\spxentry{solve\_local\_disk\_properties()}\spxextra{disk\_model.DiskModel method}}

\begin{fulllineitems}
\phantomsection\label{\detokenize{x22:disk_model.DiskModel.solve_local_disk_properties}}
\pysigstartsignatures
\pysiglinewithargsret
{\sphinxbfcode{\sphinxupquote{solve\_local\_disk\_properties}}}
{\sphinxparam{\DUrole{n}{Omega}}\sphinxparamcomma \sphinxparam{\DUrole{n}{kappa}}}
{}
\pysigstopsignatures
\end{fulllineitems}


\end{fulllineitems}

\index{DiskImage (class in disk\_model)@\spxentry{DiskImage}\spxextra{class in disk\_model}}

\begin{fulllineitems}
\phantomsection\label{\detokenize{x22:disk_model.DiskImage}}
\pysigstartsignatures
\pysiglinewithargsret
{\sphinxbfcode{\sphinxupquote{\DUrole{k}{class}\DUrole{w}{ }}}\sphinxcode{\sphinxupquote{disk\_model.}}\sphinxbfcode{\sphinxupquote{DiskImage}}}
{\sphinxparam{\DUrole{n}{fname}}\sphinxparamcomma \sphinxparam{\DUrole{n}{ra\_deg}}\sphinxparamcomma \sphinxparam{\DUrole{n}{dec\_deg}}\sphinxparamcomma \sphinxparam{\DUrole{n}{distance\_pc}}\sphinxparamcomma \sphinxparam{\DUrole{n}{rms\_Jy}}\sphinxparamcomma \sphinxparam{\DUrole{n}{disk\_pa}}\sphinxparamcomma \sphinxparam{\DUrole{n}{img\_size\_au}\DUrole{o}{=}\DUrole{default_value}{400}}\sphinxparamcomma \sphinxparam{\DUrole{n}{remove\_background}\DUrole{o}{=}\DUrole{default_value}{False}}}
{}
\pysigstopsignatures
\sphinxAtStartPar
Bases: \sphinxcode{\sphinxupquote{object}}

\sphinxAtStartPar
Stores the image of a system (at one wavelength).
Can also be used to generate mock observation for given F(R) and
compare mock observation with image.
\index{img (disk\_model.DiskImage attribute)@\spxentry{img}\spxextra{disk\_model.DiskImage attribute}}

\begin{fulllineitems}
\phantomsection\label{\detokenize{x22:disk_model.DiskImage.img}}
\pysigstartsignatures
\pysigline
{\sphinxbfcode{\sphinxupquote{img}}}
\pysigstopsignatures
\sphinxAtStartPar
observed image

\end{fulllineitems}

\index{img\_model (disk\_model.DiskImage attribute)@\spxentry{img\_model}\spxextra{disk\_model.DiskImage attribute}}

\begin{fulllineitems}
\phantomsection\label{\detokenize{x22:disk_model.DiskImage.img_model}}
\pysigstartsignatures
\pysigline
{\sphinxbfcode{\sphinxupquote{img\_model}}}
\pysigstopsignatures
\sphinxAtStartPar
mock observation image

\end{fulllineitems}

\index{au\_per\_pix (disk\_model.DiskImage attribute)@\spxentry{au\_per\_pix}\spxextra{disk\_model.DiskImage attribute}}

\begin{fulllineitems}
\phantomsection\label{\detokenize{x22:disk_model.DiskImage.au_per_pix}}
\pysigstartsignatures
\pysigline
{\sphinxbfcode{\sphinxupquote{au\_per\_pix}}}
\pysigstopsignatures
\sphinxAtStartPar
au per pixel

\end{fulllineitems}



\begin{fulllineitems}

\pysigstartsignatures
\pysigline
{\sphinxbfcode{\sphinxupquote{(see~others~in~\_\_init\_\_)}}}
\pysigstopsignatures
\end{fulllineitems}

\subsubsection*{Methods}


\begin{savenotes}
\sphinxatlongtablestart
\sphinxthistablewithglobalstyle
\sphinxthistablewithnovlinesstyle
\makeatletter
  \LTleft \@totalleftmargin plus1fill
  \LTright\dimexpr\columnwidth-\@totalleftmargin-\linewidth\relax plus1fill
\makeatother
\begin{longtable}{\X{1}{2}\X{1}{2}}
\sphinxtoprule
\endfirsthead

\multicolumn{2}{c}{\sphinxnorowcolor
    \makebox[0pt]{\sphinxtablecontinued{\tablename\ \thetable{} \textendash{} continued from previous page}}%
}\\
\sphinxtoprule
\endhead

\sphinxbottomrule
\multicolumn{2}{r}{\sphinxnorowcolor
    \makebox[0pt][r]{\sphinxtablecontinued{continues on next page}}%
}\\
\endfoot

\endlastfoot
\sphinxtableatstartofbodyhook

\sphinxAtStartPar
{\hyperref[\detokenize{x22:disk_model.DiskImage.generate_mask}]{\sphinxcrossref{\sphinxcode{\sphinxupquote{generate\_mask}}}}}(mask\_radius, cosI)
&
\sphinxAtStartPar
Generate a mask covering radius (in au) between mask\_radius.
\\
\sphinxhline
\sphinxAtStartPar
{\hyperref[\detokenize{x22:disk_model.DiskImage.generate_mock_observation}]{\sphinxcrossref{\sphinxcode{\sphinxupquote{generate\_mock\_observation}}}}}(R, I, cosI)
&
\sphinxAtStartPar
Generate mock observation for gievn I(R).
\\
\sphinxbottomrule
\end{longtable}
\sphinxtableafterendhook
\sphinxatlongtableend
\end{savenotes}


\begin{savenotes}\sphinxattablestart
\sphinxthistablewithglobalstyle
\centering
\begin{tabulary}{\linewidth}[t]{TT}
\sphinxtoprule
\sphinxtableatstartofbodyhook
\sphinxAtStartPar
\sphinxstylestrong{evaluate\_disk\_chi\_sq}
&\\
\sphinxhline
\sphinxAtStartPar
\sphinxstylestrong{evaluate\_log\_likelihood}
&\\
\sphinxhline
\sphinxAtStartPar
\sphinxstylestrong{evaluate\_log\_likelihood\_old}
&\\
\sphinxbottomrule
\end{tabulary}
\sphinxtableafterendhook\par
\sphinxattableend\end{savenotes}
\index{evaluate\_disk\_chi\_sq() (disk\_model.DiskImage method)@\spxentry{evaluate\_disk\_chi\_sq()}\spxextra{disk\_model.DiskImage method}}

\begin{fulllineitems}
\phantomsection\label{\detokenize{x22:disk_model.DiskImage.evaluate_disk_chi_sq}}
\pysigstartsignatures
\pysiglinewithargsret
{\sphinxbfcode{\sphinxupquote{evaluate\_disk\_chi\_sq}}}
{\sphinxparam{\DUrole{n}{sigma\_log\_model}\DUrole{o}{=}\DUrole{default_value}{0.34657359027997264}}}
{}
\pysigstopsignatures
\end{fulllineitems}

\index{evaluate\_log\_likelihood() (disk\_model.DiskImage method)@\spxentry{evaluate\_log\_likelihood()}\spxextra{disk\_model.DiskImage method}}

\begin{fulllineitems}
\phantomsection\label{\detokenize{x22:disk_model.DiskImage.evaluate_log_likelihood}}
\pysigstartsignatures
\pysiglinewithargsret
{\sphinxbfcode{\sphinxupquote{evaluate\_log\_likelihood}}}
{\sphinxparam{\DUrole{n}{sigma\_log\_model}\DUrole{o}{=}\DUrole{default_value}{0.34657359027997264}}}
{}
\pysigstopsignatures
\end{fulllineitems}

\index{evaluate\_log\_likelihood\_old() (disk\_model.DiskImage method)@\spxentry{evaluate\_log\_likelihood\_old()}\spxextra{disk\_model.DiskImage method}}

\begin{fulllineitems}
\phantomsection\label{\detokenize{x22:disk_model.DiskImage.evaluate_log_likelihood_old}}
\pysigstartsignatures
\pysiglinewithargsret
{\sphinxbfcode{\sphinxupquote{evaluate\_log\_likelihood\_old}}}
{\sphinxparam{\DUrole{n}{sigma\_log\_model}\DUrole{o}{=}\DUrole{default_value}{0.34657359027997264}}}
{}
\pysigstopsignatures
\end{fulllineitems}

\index{generate\_mask() (disk\_model.DiskImage method)@\spxentry{generate\_mask()}\spxextra{disk\_model.DiskImage method}}

\begin{fulllineitems}
\phantomsection\label{\detokenize{x22:disk_model.DiskImage.generate_mask}}
\pysigstartsignatures
\pysiglinewithargsret
{\sphinxbfcode{\sphinxupquote{generate\_mask}}}
{\sphinxparam{\DUrole{n}{mask\_radius}}\sphinxparamcomma \sphinxparam{\DUrole{n}{cosI}}}
{}
\pysigstopsignatures
\sphinxAtStartPar
Generate a mask covering radius (in au) between mask\_radius.
\begin{quote}\begin{description}
\sphinxlineitem{Parameters}\begin{itemize}
\item {} 
\sphinxAtStartPar
\sphinxstyleliteralstrong{\sphinxupquote{mask\_radius}} \textendash{} (2,) array

\item {} 
\sphinxAtStartPar
\sphinxstyleliteralstrong{\sphinxupquote{cosI}} \textendash{} disk inclination

\end{itemize}

\sphinxlineitem{Returns}
\sphinxAtStartPar
2d array of mask

\sphinxlineitem{Return type}
\sphinxAtStartPar
mask

\end{description}\end{quote}

\end{fulllineitems}

\index{generate\_mock\_observation() (disk\_model.DiskImage method)@\spxentry{generate\_mock\_observation()}\spxextra{disk\_model.DiskImage method}}

\begin{fulllineitems}
\phantomsection\label{\detokenize{x22:disk_model.DiskImage.generate_mock_observation}}
\pysigstartsignatures
\pysiglinewithargsret
{\sphinxbfcode{\sphinxupquote{generate\_mock\_observation}}}
{\sphinxparam{\DUrole{n}{R}}\sphinxparamcomma \sphinxparam{\DUrole{n}{I}}\sphinxparamcomma \sphinxparam{\DUrole{n}{cosI}}}
{}
\pysigstopsignatures
\sphinxAtStartPar
Generate mock observation for gievn I(R). (here I is the intensity)
\begin{quote}\begin{description}
\sphinxlineitem{Parameters}\begin{itemize}
\item {} 
\sphinxAtStartPar
\sphinxstyleliteralstrong{\sphinxupquote{R}} \textendash{} 1d array, I = I(R) at R{[}1:{]}
so len(F) = len(R)\sphinxhyphen{}1

\item {} 
\sphinxAtStartPar
\sphinxstyleliteralstrong{\sphinxupquote{I}} \textendash{} 1d array, I = I(R) at R{[}1:{]}
so len(F) = len(R)\sphinxhyphen{}1

\item {} 
\sphinxAtStartPar
\sphinxstyleliteralstrong{\sphinxupquote{cosI}} \textendash{} inclination

\end{itemize}

\end{description}\end{quote}

\end{fulllineitems}


\end{fulllineitems}

\index{DiskFitting (class in disk\_model)@\spxentry{DiskFitting}\spxextra{class in disk\_model}}

\begin{fulllineitems}
\phantomsection\label{\detokenize{x22:disk_model.DiskFitting}}
\pysigstartsignatures
\pysiglinewithargsret
{\sphinxbfcode{\sphinxupquote{\DUrole{k}{class}\DUrole{w}{ }}}\sphinxcode{\sphinxupquote{disk\_model.}}\sphinxbfcode{\sphinxupquote{DiskFitting}}}
{\sphinxparam{\DUrole{n}{source\_name}}\sphinxparamcomma \sphinxparam{\DUrole{n}{opacity\_table}}\sphinxparamcomma \sphinxparam{\DUrole{n}{disk\_property\_table}}}
{}
\pysigstopsignatures
\sphinxAtStartPar
Bases: \sphinxcode{\sphinxupquote{object}}

\sphinxAtStartPar
Class for fitting multi\sphinxhyphen{}wavelength observations of a system.
\index{source\_name (disk\_model.DiskFitting attribute)@\spxentry{source\_name}\spxextra{disk\_model.DiskFitting attribute}}

\begin{fulllineitems}
\phantomsection\label{\detokenize{x22:disk_model.DiskFitting.source_name}}
\pysigstartsignatures
\pysigline
{\sphinxbfcode{\sphinxupquote{source\_name}}}
\pysigstopsignatures
\end{fulllineitems}

\index{disk\_model (disk\_model.DiskFitting attribute)@\spxentry{disk\_model}\spxextra{disk\_model.DiskFitting attribute}}

\begin{fulllineitems}
\phantomsection\label{\detokenize{x22:disk_model.DiskFitting.disk_model}}
\pysigstartsignatures
\pysigline
{\sphinxbfcode{\sphinxupquote{disk\_model}}}
\pysigstopsignatures
\sphinxAtStartPar
DiskModel object

\end{fulllineitems}

\index{disk\_image\_list (disk\_model.DiskFitting attribute)@\spxentry{disk\_image\_list}\spxextra{disk\_model.DiskFitting attribute}}

\begin{fulllineitems}
\phantomsection\label{\detokenize{x22:disk_model.DiskFitting.disk_image_list}}
\pysigstartsignatures
\pysigline
{\sphinxbfcode{\sphinxupquote{disk\_image\_list}}}
\pysigstopsignatures
\sphinxAtStartPar
list of DiskImage objects

\end{fulllineitems}

\subsubsection*{Methods}


\begin{savenotes}
\sphinxatlongtablestart
\sphinxthistablewithglobalstyle
\sphinxthistablewithnovlinesstyle
\makeatletter
  \LTleft \@totalleftmargin plus1fill
  \LTright\dimexpr\columnwidth-\@totalleftmargin-\linewidth\relax plus1fill
\makeatother
\begin{longtable}{\X{1}{2}\X{1}{2}}
\sphinxtoprule
\endfirsthead

\multicolumn{2}{c}{\sphinxnorowcolor
    \makebox[0pt]{\sphinxtablecontinued{\tablename\ \thetable{} \textendash{} continued from previous page}}%
}\\
\sphinxtoprule
\endhead

\sphinxbottomrule
\multicolumn{2}{r}{\sphinxnorowcolor
    \makebox[0pt][r]{\sphinxtablecontinued{continues on next page}}%
}\\
\endfoot

\endlastfoot
\sphinxtableatstartofbodyhook

\sphinxAtStartPar
{\hyperref[\detokenize{x22:disk_model.DiskFitting.fit}]{\sphinxcrossref{\sphinxcode{\sphinxupquote{fit}}}}}({[}Mdot\_fn, Mdot\_arg, Q, weights, ...{]})
&
\sphinxAtStartPar
Fit model to observation.
\\
\sphinxbottomrule
\end{longtable}
\sphinxtableafterendhook
\sphinxatlongtableend
\end{savenotes}


\begin{savenotes}\sphinxattablestart
\sphinxthistablewithglobalstyle
\centering
\begin{tabulary}{\linewidth}[t]{TT}
\sphinxtoprule
\sphinxtableatstartofbodyhook
\sphinxAtStartPar
\sphinxstylestrong{add\_observation}
&\\
\sphinxhline
\sphinxAtStartPar
\sphinxstylestrong{evaluate\_log\_likelihood}
&\\
\sphinxhline
\sphinxAtStartPar
\sphinxstylestrong{mcmc\_fit}
&\\
\sphinxhline
\sphinxAtStartPar
\sphinxstylestrong{set\_cosI}
&\\
\sphinxbottomrule
\end{tabulary}
\sphinxtableafterendhook\par
\sphinxattableend\end{savenotes}
\index{add\_observation() (disk\_model.DiskFitting method)@\spxentry{add\_observation()}\spxextra{disk\_model.DiskFitting method}}

\begin{fulllineitems}
\phantomsection\label{\detokenize{x22:disk_model.DiskFitting.add_observation}}
\pysigstartsignatures
\pysiglinewithargsret
{\sphinxbfcode{\sphinxupquote{add\_observation}}}
{\sphinxparam{\DUrole{n}{disk\_image}}\sphinxparamcomma \sphinxparam{\DUrole{n}{lam\_obs}}}
{}
\pysigstopsignatures
\end{fulllineitems}

\index{evaluate\_log\_likelihood() (disk\_model.DiskFitting method)@\spxentry{evaluate\_log\_likelihood()}\spxextra{disk\_model.DiskFitting method}}

\begin{fulllineitems}
\phantomsection\label{\detokenize{x22:disk_model.DiskFitting.evaluate_log_likelihood}}
\pysigstartsignatures
\pysiglinewithargsret
{\sphinxbfcode{\sphinxupquote{evaluate\_log\_likelihood}}}
{\sphinxparam{\DUrole{n}{Mstar}}\sphinxparamcomma \sphinxparam{\DUrole{n}{Rd}}\sphinxparamcomma \sphinxparam{\DUrole{n}{Mdot}}\sphinxparamcomma \sphinxparam{\DUrole{n}{Q}}\sphinxparamcomma \sphinxparam{\DUrole{n}{weights}\DUrole{o}{=}\DUrole{default_value}{None}}}
{}
\pysigstopsignatures
\end{fulllineitems}

\index{fit() (disk\_model.DiskFitting method)@\spxentry{fit()}\spxextra{disk\_model.DiskFitting method}}

\begin{fulllineitems}
\phantomsection\label{\detokenize{x22:disk_model.DiskFitting.fit}}
\pysigstartsignatures
\pysiglinewithargsret
{\sphinxbfcode{\sphinxupquote{fit}}}
{\sphinxparam{\DUrole{n}{Mdot\_fn=\textless{}function const\_Mdot\_over\_Mstar\textgreater{}}}\sphinxparamcomma \sphinxparam{\DUrole{n}{Mdot\_arg=3.170979198376459e\sphinxhyphen{}13}}\sphinxparamcomma \sphinxparam{\DUrole{n}{Q=1.5}}\sphinxparamcomma \sphinxparam{\DUrole{n}{weights=None}}\sphinxparamcomma \sphinxparam{\DUrole{n}{kwargs\_for\_minimize=None}}}
{}
\pysigstopsignatures
\sphinxAtStartPar
Fit model to observation.
\begin{quote}\begin{description}
\sphinxlineitem{Parameters}\begin{itemize}
\item {} 
\sphinxAtStartPar
\sphinxstyleliteralstrong{\sphinxupquote{Mdot\_fn}} \textendash{} a function that takes (Mstar, Rd, Mdot\_arg) and
returns Mdot

\item {} 
\sphinxAtStartPar
\sphinxstyleliteralstrong{\sphinxupquote{Mdot\_arg}} \textendash{} the last argument for Mdot\_fn()

\item {} 
\sphinxAtStartPar
\sphinxstyleliteralstrong{\sphinxupquote{Q}} \textendash{} Toomre Q to be passed to disk model

\item {} 
\sphinxAtStartPar
\sphinxstyleliteralstrong{\sphinxupquote{weights}} \textendash{} weights for each observed wavelength

\item {} 
\sphinxAtStartPar
\sphinxstyleliteralstrong{\sphinxupquote{kwargs\_for\_minimize}} \textendash{} kwargs to be passed to
optimize.minimize();
see self.default\_kwargs\_for\_minimize

\end{itemize}

\end{description}\end{quote}

\end{fulllineitems}

\index{mcmc\_fit() (disk\_model.DiskFitting method)@\spxentry{mcmc\_fit()}\spxextra{disk\_model.DiskFitting method}}

\begin{fulllineitems}
\phantomsection\label{\detokenize{x22:disk_model.DiskFitting.mcmc_fit}}
\pysigstartsignatures
\pysiglinewithargsret
{\sphinxbfcode{\sphinxupquote{mcmc\_fit}}}
{\sphinxparam{\DUrole{n}{Mstar}}\sphinxparamcomma \sphinxparam{\DUrole{n}{Mdot}}}
{}
\pysigstopsignatures
\end{fulllineitems}

\index{set\_cosI() (disk\_model.DiskFitting method)@\spxentry{set\_cosI()}\spxextra{disk\_model.DiskFitting method}}

\begin{fulllineitems}
\phantomsection\label{\detokenize{x22:disk_model.DiskFitting.set_cosI}}
\pysigstartsignatures
\pysiglinewithargsret
{\sphinxbfcode{\sphinxupquote{set\_cosI}}}
{\sphinxparam{\DUrole{n}{cosI}}}
{}
\pysigstopsignatures
\end{fulllineitems}


\end{fulllineitems}


\sphinxstepscope


\chapter{Vertical Structure}
\label{\detokenize{vertical:vertical-structure}}\label{\detokenize{vertical::doc}}

\section{Formulations}
\label{\detokenize{vertical:formulations}}

\section{Grid Interpolation}
\label{\detokenize{vertical:grid-interpolation}}

\section{Modules}
\label{\detokenize{vertical:modules}}\index{DiskModel\_spherical (class in spherical\_x22)@\spxentry{DiskModel\_spherical}\spxextra{class in spherical\_x22}}

\begin{fulllineitems}
\phantomsection\label{\detokenize{vertical:spherical_x22.DiskModel_spherical}}
\pysigstartsignatures
\pysiglinewithargsret
{\sphinxbfcode{\sphinxupquote{\DUrole{k}{class}\DUrole{w}{ }}}\sphinxcode{\sphinxupquote{spherical\_x22.}}\sphinxbfcode{\sphinxupquote{DiskModel\_spherical}}}
{\sphinxparam{\DUrole{n}{opacity\_table}}\sphinxparamcomma \sphinxparam{\DUrole{n}{disk\_property\_table}}}
{}
\pysigstopsignatures
\sphinxAtStartPar
Bases: \sphinxcode{\sphinxupquote{object}}
\subsubsection*{Methods}


\begin{savenotes}
\sphinxatlongtablestart
\sphinxthistablewithglobalstyle
\sphinxthistablewithnovlinesstyle
\makeatletter
  \LTleft \@totalleftmargin plus1fill
  \LTright\dimexpr\columnwidth-\@totalleftmargin-\linewidth\relax plus1fill
\makeatother
\begin{longtable}{\X{1}{2}\X{1}{2}}
\sphinxtoprule
\endfirsthead

\multicolumn{2}{c}{\sphinxnorowcolor
    \makebox[0pt]{\sphinxtablecontinued{\tablename\ \thetable{} \textendash{} continued from previous page}}%
}\\
\sphinxtoprule
\endhead

\sphinxbottomrule
\multicolumn{2}{r}{\sphinxnorowcolor
    \makebox[0pt][r]{\sphinxtablecontinued{continues on next page}}%
}\\
\endfoot

\endlastfoot
\sphinxtableatstartofbodyhook

\sphinxAtStartPar
{\hyperref[\detokenize{vertical:spherical_x22.DiskModel_spherical.extend_to_spherical}]{\sphinxcrossref{\sphinxcode{\sphinxupquote{extend\_to\_spherical}}}}}(NTheta, NPhi{[}, ...{]})
&
\sphinxAtStartPar
Extend the disk model to spherical coordinates
\\
\sphinxhline
\sphinxAtStartPar
{\hyperref[\detokenize{vertical:spherical_x22.DiskModel_spherical.input_disk_parameter}]{\sphinxcrossref{\sphinxcode{\sphinxupquote{input\_disk\_parameter}}}}}(Mstar, Mdot, Rd, Q{[}, N\_R{]})
&
\sphinxAtStartPar
Run Wenrui\textquotesingle{}s radial profile to get initial r\sphinxhyphen{}dependent profiles.
\\
\sphinxhline
\sphinxAtStartPar
{\hyperref[\detokenize{vertical:spherical_x22.DiskModel_spherical.make_position_map}]{\sphinxcrossref{\sphinxcode{\sphinxupquote{make\_position\_map}}}}}()
&
\sphinxAtStartPar
Make grid of position in cylindrical coordinates
\\
\sphinxhline
\sphinxAtStartPar
{\hyperref[\detokenize{vertical:spherical_x22.DiskModel_spherical.make_rho_and_m_map}]{\sphinxcrossref{\sphinxcode{\sphinxupquote{make\_rho\_and\_m\_map}}}}}()
&
\sphinxAtStartPar
Calculate volume density and mass\sphinxhyphen{}depth scale
\\
\sphinxhline
\sphinxAtStartPar
{\hyperref[\detokenize{vertical:spherical_x22.DiskModel_spherical.make_spherical_grid}]{\sphinxcrossref{\sphinxcode{\sphinxupquote{make\_spherical\_grid}}}}}()
&
\sphinxAtStartPar
Making the grid which RADMC3D adopts
\\
\sphinxhline
\sphinxAtStartPar
{\hyperref[\detokenize{vertical:spherical_x22.DiskModel_spherical.make_tau_and_T_map}]{\sphinxcrossref{\sphinxcode{\sphinxupquote{make\_tau\_and\_T\_map}}}}}()
&
\sphinxAtStartPar
Using disk\_property\_table to calculate kappa\_r from T, and further calculating tau\_r and T
\\
\sphinxhline
\sphinxAtStartPar
{\hyperref[\detokenize{vertical:spherical_x22.DiskModel_spherical.precompute_property}]{\sphinxcrossref{\sphinxcode{\sphinxupquote{precompute\_property}}}}}({[}miu, ionized\_factor{]})
&
\sphinxAtStartPar
To calculate the gas pressure scale height miu   : mean molecular weight (Default = 2.3) ionized\_factor: N/(N\sphinxhyphen{}n\_e) (ionized factor; Default = 1)
\\
\sphinxbottomrule
\end{longtable}
\sphinxtableafterendhook
\sphinxatlongtableend
\end{savenotes}
\index{extend\_to\_spherical() (spherical\_x22.DiskModel\_spherical method)@\spxentry{extend\_to\_spherical()}\spxextra{spherical\_x22.DiskModel\_spherical method}}

\begin{fulllineitems}
\phantomsection\label{\detokenize{vertical:spherical_x22.DiskModel_spherical.extend_to_spherical}}
\pysigstartsignatures
\pysiglinewithargsret
{\sphinxbfcode{\sphinxupquote{extend\_to\_spherical}}}
{\sphinxparam{\DUrole{n}{NTheta}}\sphinxparamcomma \sphinxparam{\DUrole{n}{NPhi}}\sphinxparamcomma \sphinxparam{\DUrole{n}{theta\_min\_deg}\DUrole{o}{=}\DUrole{default_value}{30}}}
{}
\pysigstopsignatures
\sphinxAtStartPar
Extend the disk model to spherical coordinates
\begin{quote}\begin{description}
\sphinxlineitem{Parameters}\begin{itemize}
\item {} 
\sphinxAtStartPar
\sphinxstyleliteralstrong{\sphinxupquote{NTheta}} (\sphinxstyleliteralemphasis{\sphinxupquote{float}}) \textendash{} The number of theta grid

\item {} 
\sphinxAtStartPar
\sphinxstyleliteralstrong{\sphinxupquote{NPhi}} (\sphinxstyleliteralemphasis{\sphinxupquote{float}}) \textendash{} The number of phi grid

\item {} 
\sphinxAtStartPar
\sphinxstyleliteralstrong{\sphinxupquote{theta\_min\_deg}} (\sphinxstyleliteralemphasis{\sphinxupquote{float}}) \textendash{} The minimum angle of theta grid (Default = 30 degree)

\end{itemize}

\end{description}\end{quote}

\end{fulllineitems}

\index{input\_disk\_parameter() (spherical\_x22.DiskModel\_spherical method)@\spxentry{input\_disk\_parameter()}\spxextra{spherical\_x22.DiskModel\_spherical method}}

\begin{fulllineitems}
\phantomsection\label{\detokenize{vertical:spherical_x22.DiskModel_spherical.input_disk_parameter}}
\pysigstartsignatures
\pysiglinewithargsret
{\sphinxbfcode{\sphinxupquote{input\_disk\_parameter}}}
{\sphinxparam{\DUrole{n}{Mstar}}\sphinxparamcomma \sphinxparam{\DUrole{n}{Mdot}}\sphinxparamcomma \sphinxparam{\DUrole{n}{Rd}}\sphinxparamcomma \sphinxparam{\DUrole{n}{Q}}\sphinxparamcomma \sphinxparam{\DUrole{n}{N\_R}\DUrole{o}{=}\DUrole{default_value}{500}}}
{}
\pysigstopsignatures
\sphinxAtStartPar
Run Wenrui’s radial profile to get initial r\sphinxhyphen{}dependent profiles.
\begin{quote}\begin{description}
\sphinxlineitem{Parameters}\begin{itemize}
\item {} 
\sphinxAtStartPar
\sphinxstyleliteralstrong{\sphinxupquote{Mstar}} (\sphinxstyleliteralemphasis{\sphinxupquote{float}}) \textendash{} Mass of protostar

\item {} 
\sphinxAtStartPar
\sphinxstyleliteralstrong{\sphinxupquote{Mdot}} (\sphinxstyleliteralemphasis{\sphinxupquote{float}}) \textendash{} Rate of mass infall from the envelope onto the disk

\item {} 
\sphinxAtStartPar
\sphinxstyleliteralstrong{\sphinxupquote{Rd}} (\sphinxstyleliteralemphasis{\sphinxupquote{float}}) \textendash{} Radius of the disk

\item {} 
\sphinxAtStartPar
\sphinxstyleliteralstrong{\sphinxupquote{Q}} (\sphinxstyleliteralemphasis{\sphinxupquote{float}}) \textendash{} Toomre index

\item {} 
\sphinxAtStartPar
\sphinxstyleliteralstrong{\sphinxupquote{N\_R}} (\sphinxstyleliteralemphasis{\sphinxupquote{float}}) \textendash{} Resolution of radius grid (Default = 500)

\end{itemize}

\end{description}\end{quote}

\end{fulllineitems}

\index{make\_position\_map() (spherical\_x22.DiskModel\_spherical method)@\spxentry{make\_position\_map()}\spxextra{spherical\_x22.DiskModel\_spherical method}}

\begin{fulllineitems}
\phantomsection\label{\detokenize{vertical:spherical_x22.DiskModel_spherical.make_position_map}}
\pysigstartsignatures
\pysiglinewithargsret
{\sphinxbfcode{\sphinxupquote{make\_position\_map}}}
{}
{}
\pysigstopsignatures
\sphinxAtStartPar
Make grid of position in cylindrical coordinates

\end{fulllineitems}

\index{make\_rho\_and\_m\_map() (spherical\_x22.DiskModel\_spherical method)@\spxentry{make\_rho\_and\_m\_map()}\spxextra{spherical\_x22.DiskModel\_spherical method}}

\begin{fulllineitems}
\phantomsection\label{\detokenize{vertical:spherical_x22.DiskModel_spherical.make_rho_and_m_map}}
\pysigstartsignatures
\pysiglinewithargsret
{\sphinxbfcode{\sphinxupquote{make\_rho\_and\_m\_map}}}
{}
{}
\pysigstopsignatures
\sphinxAtStartPar
Calculate volume density and mass\sphinxhyphen{}depth scale

\end{fulllineitems}

\index{make\_spherical\_grid() (spherical\_x22.DiskModel\_spherical method)@\spxentry{make\_spherical\_grid()}\spxextra{spherical\_x22.DiskModel\_spherical method}}

\begin{fulllineitems}
\phantomsection\label{\detokenize{vertical:spherical_x22.DiskModel_spherical.make_spherical_grid}}
\pysigstartsignatures
\pysiglinewithargsret
{\sphinxbfcode{\sphinxupquote{make\_spherical\_grid}}}
{}
{}
\pysigstopsignatures
\sphinxAtStartPar
Making the grid which RADMC3D adopts

\end{fulllineitems}

\index{make\_tau\_and\_T\_map() (spherical\_x22.DiskModel\_spherical method)@\spxentry{make\_tau\_and\_T\_map()}\spxextra{spherical\_x22.DiskModel\_spherical method}}

\begin{fulllineitems}
\phantomsection\label{\detokenize{vertical:spherical_x22.DiskModel_spherical.make_tau_and_T_map}}
\pysigstartsignatures
\pysiglinewithargsret
{\sphinxbfcode{\sphinxupquote{make\_tau\_and\_T\_map}}}
{}
{}
\pysigstopsignatures
\sphinxAtStartPar
Using disk\_property\_table to calculate kappa\_r from T, and further calculating tau\_r and T

\end{fulllineitems}

\index{precompute\_property() (spherical\_x22.DiskModel\_spherical method)@\spxentry{precompute\_property()}\spxextra{spherical\_x22.DiskModel\_spherical method}}

\begin{fulllineitems}
\phantomsection\label{\detokenize{vertical:spherical_x22.DiskModel_spherical.precompute_property}}
\pysigstartsignatures
\pysiglinewithargsret
{\sphinxbfcode{\sphinxupquote{precompute\_property}}}
{\sphinxparam{\DUrole{n}{miu}\DUrole{o}{=}\DUrole{default_value}{2.3}}\sphinxparamcomma \sphinxparam{\DUrole{n}{ionized\_factor}\DUrole{o}{=}\DUrole{default_value}{1}}}
{}
\pysigstopsignatures
\sphinxAtStartPar
To calculate the gas pressure scale height
miu   : mean molecular weight (Default = 2.3)
ionized\_factor: N/(N\sphinxhyphen{}n\_e) (ionized factor; Default = 1)

\end{fulllineitems}


\end{fulllineitems}


\sphinxstepscope


\chapter{Pipelines to RADMC\sphinxhyphen{}3D}
\label{\detokenize{radmc3d:pipelines-to-radmc-3d}}\label{\detokenize{radmc3d::doc}}

\section{Setup}
\label{\detokenize{radmc3d:setup}}

\section{Simulations}
\label{\detokenize{radmc3d:simulations}}

\subsection{Dust Continuum}
\label{\detokenize{radmc3d:dust-continuum}}

\subsection{Molecular Line Emissions}
\label{\detokenize{radmc3d:molecular-line-emissions}}

\subsection{Spectral Energy Distribution (SED)}
\label{\detokenize{radmc3d:spectral-energy-distribution-sed}}

\section{Modules}
\label{\detokenize{radmc3d:modules}}\index{radmc3d\_setup (class in setup)@\spxentry{radmc3d\_setup}\spxextra{class in setup}}

\begin{fulllineitems}
\phantomsection\label{\detokenize{radmc3d:setup.radmc3d_setup}}
\pysigstartsignatures
\pysiglinewithargsret
{\sphinxbfcode{\sphinxupquote{\DUrole{k}{class}\DUrole{w}{ }}}\sphinxcode{\sphinxupquote{setup.}}\sphinxbfcode{\sphinxupquote{radmc3d\_setup}}}
{\sphinxparam{\DUrole{n}{silent}\DUrole{o}{=}\DUrole{default_value}{True}}}
{}
\pysigstopsignatures
\sphinxAtStartPar
Bases: \sphinxcode{\sphinxupquote{object}}

\sphinxAtStartPar
A class to initialize the radmc3d simulation.
\subsubsection*{Examples}

\sphinxAtStartPar
Simplest usage (defaulting everything):

\begin{sphinxVerbatim}[commandchars=\\\{\}]
\PYG{n}{test} \PYG{o}{=} \PYG{n}{radmc3d\PYGZus{}setup}\PYG{p}{(}\PYG{p}{)}
\PYG{n}{test}\PYG{o}{.}\PYG{n}{get\PYGZus{}mastercontrol}\PYG{p}{(}\PYG{p}{)}
\PYG{n}{test}\PYG{o}{.}\PYG{n}{get\PYGZus{}continuumlambda}\PYG{p}{(}\PYG{p}{)}
\end{sphinxVerbatim}

\sphinxAtStartPar
More complicated cases:

\begin{sphinxVerbatim}[commandchars=\\\{\}]
\PYG{n}{test} \PYG{o}{=} \PYG{n}{radmc3d\PYGZus{}setup}\PYG{p}{(}\PYG{n}{silent}\PYG{o}{=}\PYG{k+kc}{False}\PYG{p}{)}

\PYG{n}{test}\PYG{o}{.}\PYG{n}{get\PYGZus{}mastercontrol}\PYG{p}{(}\PYG{n}{filename}\PYG{o}{=}\PYG{l+s+s1}{\PYGZsq{}}\PYG{l+s+s1}{radmctest.inp}\PYG{l+s+s1}{\PYGZsq{}}\PYG{p}{,}
                       \PYG{n}{comment}\PYG{o}{=}\PYG{l+s+s1}{\PYGZsq{}}\PYG{l+s+s1}{this is a test}\PYG{l+s+s1}{\PYGZsq{}}\PYG{p}{,}
                       \PYG{n}{incl\PYGZus{}dust}\PYG{o}{=}\PYG{l+m+mi}{1}\PYG{p}{)}

\PYG{n}{test}\PYG{o}{.}\PYG{n}{get\PYGZus{}linecontrol}\PYG{p}{(}\PYG{n}{filename}\PYG{o}{=}\PYG{l+s+s1}{\PYGZsq{}}\PYG{l+s+s1}{lines\PYGZus{}test.inp}\PYG{l+s+s1}{\PYGZsq{}}\PYG{p}{,}
                     \PYG{n}{comment}\PYG{o}{=}\PYG{l+s+s1}{\PYGZsq{}}\PYG{l+s+s1}{this is a test}\PYG{l+s+s1}{\PYGZsq{}}\PYG{p}{,}
                     \PYG{n}{line1}\PYG{o}{=}\PYG{l+s+s1}{\PYGZsq{}}\PYG{l+s+s1}{ch3oh leiden 0 0}\PYG{l+s+s1}{\PYGZsq{}}\PYG{p}{,}
                     \PYG{n}{line2}\PYG{o}{=}\PYG{l+s+s1}{\PYGZsq{}}\PYG{l+s+s1}{co leiden 0 0}\PYG{l+s+s1}{\PYGZsq{}}\PYG{p}{)}

\PYG{n}{lam1}\PYG{p}{,} \PYG{n}{lam2}\PYG{p}{,} \PYG{n}{lam3}\PYG{p}{,} \PYG{n}{lam4} \PYG{o}{=} \PYG{l+m+mf}{0.1e0}\PYG{p}{,} \PYG{l+m+mf}{1.0e2}\PYG{p}{,} \PYG{l+m+mf}{5.0e3}\PYG{p}{,} \PYG{l+m+mf}{1.0e4}
\PYG{n}{n12}\PYG{p}{,} \PYG{n}{n23}\PYG{p}{,} \PYG{n}{n34} \PYG{o}{=} \PYG{l+m+mi}{100}\PYG{p}{,} \PYG{l+m+mi}{100}\PYG{p}{,} \PYG{l+m+mi}{50}
\PYG{n}{lam12} \PYG{o}{=} \PYG{n}{np}\PYG{o}{.}\PYG{n}{logspace}\PYG{p}{(}\PYG{n}{np}\PYG{o}{.}\PYG{n}{log10}\PYG{p}{(}\PYG{n}{lam1}\PYG{p}{)}\PYG{p}{,} \PYG{n}{np}\PYG{o}{.}\PYG{n}{log10}\PYG{p}{(}\PYG{n}{lam2}\PYG{p}{)}\PYG{p}{,} \PYG{n}{n12}\PYG{p}{,} \PYG{n}{endpoint}\PYG{o}{=}\PYG{k+kc}{False}\PYG{p}{)}
\PYG{n}{lam23} \PYG{o}{=} \PYG{n}{np}\PYG{o}{.}\PYG{n}{logspace}\PYG{p}{(}\PYG{n}{np}\PYG{o}{.}\PYG{n}{log10}\PYG{p}{(}\PYG{n}{lam2}\PYG{p}{)}\PYG{p}{,} \PYG{n}{np}\PYG{o}{.}\PYG{n}{log10}\PYG{p}{(}\PYG{n}{lam3}\PYG{p}{)}\PYG{p}{,} \PYG{n}{n23}\PYG{p}{,} \PYG{n}{endpoint}\PYG{o}{=}\PYG{k+kc}{False}\PYG{p}{)}
\PYG{n}{lam34} \PYG{o}{=} \PYG{n}{np}\PYG{o}{.}\PYG{n}{logspace}\PYG{p}{(}\PYG{n}{np}\PYG{o}{.}\PYG{n}{log10}\PYG{p}{(}\PYG{n}{lam3}\PYG{p}{)}\PYG{p}{,} \PYG{n}{np}\PYG{o}{.}\PYG{n}{log10}\PYG{p}{(}\PYG{n}{lam4}\PYG{p}{)}\PYG{p}{,} \PYG{n}{n34}\PYG{p}{,} \PYG{n}{endpoint}\PYG{o}{=}\PYG{k+kc}{True}\PYG{p}{)}
\PYG{n}{lams} \PYG{o}{=} \PYG{p}{[}\PYG{n}{lam12}\PYG{p}{,} \PYG{n}{lam23}\PYG{p}{,} \PYG{n}{lam34}\PYG{p}{]}

\PYG{k}{for} \PYG{n}{i} \PYG{o+ow}{in} \PYG{n+nb}{range}\PYG{p}{(}\PYG{n+nb}{len}\PYG{p}{(}\PYG{n}{lams}\PYG{p}{)}\PYG{p}{)}\PYG{p}{:}
    \PYG{n}{append} \PYG{o}{=} \PYG{p}{(}\PYG{n}{i} \PYG{o}{!=} \PYG{l+m+mi}{0}\PYG{p}{)}
    \PYG{n}{test}\PYG{o}{.}\PYG{n}{get\PYGZus{}continuumlambda}\PYG{p}{(}\PYG{n}{filename}\PYG{o}{=}\PYG{l+s+s1}{\PYGZsq{}}\PYG{l+s+s1}{wavelength\PYGZus{}micron\PYGZus{}test.inp}\PYG{l+s+s1}{\PYGZsq{}}\PYG{p}{,}
                             \PYG{n}{comment}\PYG{o}{=}\PYG{l+s+s1}{\PYGZsq{}}\PYG{l+s+s1}{this is a test}\PYG{l+s+s1}{\PYGZsq{}}\PYG{p}{,}
                             \PYG{n}{lambda\PYGZus{}micron}\PYG{o}{=}\PYG{n}{lams}\PYG{p}{[}\PYG{n}{i}\PYG{p}{]}\PYG{p}{,} \PYG{n}{append}\PYG{o}{=}\PYG{n}{append}\PYG{p}{)}
\end{sphinxVerbatim}
\subsubsection*{Methods}


\begin{savenotes}
\sphinxatlongtablestart
\sphinxthistablewithglobalstyle
\sphinxthistablewithnovlinesstyle
\makeatletter
  \LTleft \@totalleftmargin plus1fill
  \LTright\dimexpr\columnwidth-\@totalleftmargin-\linewidth\relax plus1fill
\makeatother
\begin{longtable}{\X{1}{2}\X{1}{2}}
\sphinxtoprule
\endfirsthead

\multicolumn{2}{c}{\sphinxnorowcolor
    \makebox[0pt]{\sphinxtablecontinued{\tablename\ \thetable{} \textendash{} continued from previous page}}%
}\\
\sphinxtoprule
\endhead

\sphinxbottomrule
\multicolumn{2}{r}{\sphinxnorowcolor
    \makebox[0pt][r]{\sphinxtablecontinued{continues on next page}}%
}\\
\endfoot

\endlastfoot
\sphinxtableatstartofbodyhook

\sphinxAtStartPar
{\hyperref[\detokenize{radmc3d:setup.radmc3d_setup.duplicate_file}]{\sphinxcrossref{\sphinxcode{\sphinxupquote{duplicate\_file}}}}}(default\_filename, filename{[}, ...{]})
&
\sphinxAtStartPar
A small function to duplicate the .inp files.
\\
\sphinxhline
\sphinxAtStartPar
{\hyperref[\detokenize{radmc3d:setup.radmc3d_setup.get_continuumlambda}]{\sphinxcrossref{\sphinxcode{\sphinxupquote{get\_continuumlambda}}}}}({[}filename, comment, ...{]})
&
\sphinxAtStartPar
Preparing the wavelength file (wavelengths are in units of micron).
\\
\sphinxhline
\sphinxAtStartPar
{\hyperref[\detokenize{radmc3d:setup.radmc3d_setup.get_diskcontrol}]{\sphinxcrossref{\sphinxcode{\sphinxupquote{get\_diskcontrol}}}}}({[}d\_to\_g\_ratio, a\_max, q, ...{]})
&
\sphinxAtStartPar
Preparing the control file for disk model.
\\
\sphinxhline
\sphinxAtStartPar
{\hyperref[\detokenize{radmc3d:setup.radmc3d_setup.get_gasdensitycontrol}]{\sphinxcrossref{\sphinxcode{\sphinxupquote{get\_gasdensitycontrol}}}}}({[}abundance, snowline, ...{]})
&
\sphinxAtStartPar
Preparing the control file for gas density.
\\
\sphinxhline
\sphinxAtStartPar
{\hyperref[\detokenize{radmc3d:setup.radmc3d_setup.get_heatcontrol}]{\sphinxcrossref{\sphinxcode{\sphinxupquote{get\_heatcontrol}}}}}({[}L\_star, R\_star, accretion, ...{]})
&
\sphinxAtStartPar
Preparing the control file for heating mechanism.
\\
\sphinxhline
\sphinxAtStartPar
{\hyperref[\detokenize{radmc3d:setup.radmc3d_setup.get_linecontrol}]{\sphinxcrossref{\sphinxcode{\sphinxupquote{get\_linecontrol}}}}}({[}filename, comment{]})
&
\sphinxAtStartPar
Preparing the control file for lines.
\\
\sphinxhline
\sphinxAtStartPar
{\hyperref[\detokenize{radmc3d:setup.radmc3d_setup.get_mastercontrol}]{\sphinxcrossref{\sphinxcode{\sphinxupquote{get\_mastercontrol}}}}}({[}filename, comment, ...{]})
&
\sphinxAtStartPar
Preparing the master control file for radmc3d.
\\
\sphinxhline
\sphinxAtStartPar
{\hyperref[\detokenize{radmc3d:setup.radmc3d_setup.get_vfieldcontrol}]{\sphinxcrossref{\sphinxcode{\sphinxupquote{get\_vfieldcontrol}}}}}({[}Kep, vinfall, Rcb, outflow{]})
&
\sphinxAtStartPar
Preparing the control file for velocity field.
\\
\sphinxhline
\sphinxAtStartPar
{\hyperref[\detokenize{radmc3d:setup.radmc3d_setup.write_dust_density}]{\sphinxcrossref{\sphinxcode{\sphinxupquote{write\_dust\_density}}}}}()
&
\sphinxAtStartPar
Preparing the control file for dust density.
\\
\sphinxhline
\sphinxAtStartPar
{\hyperref[\detokenize{radmc3d:setup.radmc3d_setup.write_dust_opac}]{\sphinxcrossref{\sphinxcode{\sphinxupquote{write\_dust\_opac}}}}}()
&
\sphinxAtStartPar
Preparing the control file for dust opacity.
\\
\sphinxbottomrule
\end{longtable}
\sphinxtableafterendhook
\sphinxatlongtableend
\end{savenotes}


\begin{savenotes}\sphinxattablestart
\sphinxthistablewithglobalstyle
\centering
\begin{tabulary}{\linewidth}[t]{TT}
\sphinxtoprule
\sphinxtableatstartofbodyhook
\sphinxAtStartPar
\sphinxstylestrong{write\_amr\_grid}
&\\
\sphinxbottomrule
\end{tabulary}
\sphinxtableafterendhook\par
\sphinxattableend\end{savenotes}
\index{duplicate\_file() (setup.radmc3d\_setup method)@\spxentry{duplicate\_file()}\spxextra{setup.radmc3d\_setup method}}

\begin{fulllineitems}
\phantomsection\label{\detokenize{radmc3d:setup.radmc3d_setup.duplicate_file}}
\pysigstartsignatures
\pysiglinewithargsret
{\sphinxbfcode{\sphinxupquote{duplicate\_file}}}
{\sphinxparam{\DUrole{n}{default\_filename}}\sphinxparamcomma \sphinxparam{\DUrole{n}{filename}}\sphinxparamcomma \sphinxparam{\DUrole{n}{comment}\DUrole{o}{=}\DUrole{default_value}{None}}\sphinxparamcomma \sphinxparam{\DUrole{n}{timemark}\DUrole{o}{=}\DUrole{default_value}{None}}}
{}
\pysigstopsignatures
\sphinxAtStartPar
A small function to duplicate the .inp files.
\begin{quote}\begin{description}
\sphinxlineitem{Parameters}\begin{itemize}
\item {} 
\sphinxAtStartPar
\sphinxstyleliteralstrong{\sphinxupquote{default\_filename}} (\sphinxstyleliteralemphasis{\sphinxupquote{str}}) \textendash{} Default output file name

\item {} 
\sphinxAtStartPar
\sphinxstyleliteralstrong{\sphinxupquote{filename}} (\sphinxstyleliteralemphasis{\sphinxupquote{str}}) \textendash{} Filename of the duplication

\item {} 
\sphinxAtStartPar
\sphinxstyleliteralstrong{\sphinxupquote{comment}} (\sphinxstyleliteralemphasis{\sphinxupquote{str}}) \textendash{} If present, include it as header in the duplicated output file.

\item {} 
\sphinxAtStartPar
\sphinxstyleliteralstrong{\sphinxupquote{timemark}} (\sphinxstyleliteralemphasis{\sphinxupquote{str}}) \textendash{} If present, include it as header in the duplicated output file.

\end{itemize}

\end{description}\end{quote}

\end{fulllineitems}

\index{get\_continuumlambda() (setup.radmc3d\_setup method)@\spxentry{get\_continuumlambda()}\spxextra{setup.radmc3d\_setup method}}

\begin{fulllineitems}
\phantomsection\label{\detokenize{radmc3d:setup.radmc3d_setup.get_continuumlambda}}
\pysigstartsignatures
\pysiglinewithargsret
{\sphinxbfcode{\sphinxupquote{get\_continuumlambda}}}
{\sphinxparam{\DUrole{n}{filename}\DUrole{o}{=}\DUrole{default_value}{None}}\sphinxparamcomma \sphinxparam{\DUrole{n}{comment}\DUrole{o}{=}\DUrole{default_value}{None}}\sphinxparamcomma \sphinxparam{\DUrole{n}{lambda\_micron}\DUrole{o}{=}\DUrole{default_value}{None}}\sphinxparamcomma \sphinxparam{\DUrole{n}{append}\DUrole{o}{=}\DUrole{default_value}{False}}\sphinxparamcomma \sphinxparam{\DUrole{n}{silent}\DUrole{o}{=}\DUrole{default_value}{False}}}
{}
\pysigstopsignatures
\sphinxAtStartPar
Preparing the wavelength file (wavelengths are in units of micron).
This is the file that sets the discrete wavelength points for the continuum radiative transfer calculations.
If this method is called for multiple times, by default it concatenate the wavelengths in each input,
unless the input wavelengths do not increase or decrease monotonically.
In that case, it gives a warning without editing the files.
If no lambda\_micron is given, it recreates the ‘wavelength\_micron.inp’ file using the default wavelengths.

\sphinxAtStartPar
Format:
nlam
lambda{[}1{]}
…
…
lambda{[}nlam{]}

\sphinxAtStartPar
\sphinxurl{https://www.ita.uni-heidelberg.de/~dullemond/software/radmc-3d/manual\_radmc3d/inputoutputfiles.html\#sec-wavelengths}
\begin{quote}\begin{description}
\sphinxlineitem{Parameters}\begin{itemize}
\item {} 
\sphinxAtStartPar
\sphinxstyleliteralstrong{\sphinxupquote{filename}} (\sphinxstyleliteralemphasis{\sphinxupquote{string}}) \textendash{} output filename. (default: wavelength\_micron.inp).
It will still creat a file with default name,
but will duplicate an output file with the specified name.

\item {} 
\sphinxAtStartPar
\sphinxstyleliteralstrong{\sphinxupquote{comment}} (\sphinxstyleliteralemphasis{\sphinxupquote{string}}) \textendash{} comment to add to the file header (default: None)

\item {} 
\sphinxAtStartPar
\sphinxstyleliteralstrong{\sphinxupquote{lambda\_micron}} (\sphinxstyleliteralemphasis{\sphinxupquote{numpy array}}) \textendash{} wavelength to calculate continuum (in units of micron).

\item {} 
\sphinxAtStartPar
\sphinxstyleliteralstrong{\sphinxupquote{append}} (\sphinxstyleliteralemphasis{\sphinxupquote{bool}}) \textendash{} if False, remove the existing wavelength\_micron.inp and ignore any information in it (default: False)

\end{itemize}

\end{description}\end{quote}

\begin{sphinxadmonition}{note}{Note:}
\sphinxAtStartPar
Wavelengths must be monotonically increasing/decreasing.
\end{sphinxadmonition}

\begin{sphinxadmonition}{note}{Note:}
\sphinxAtStartPar
Wavelength coverage must include the wavelengths at which the stellar spectra have most of their energy,
and at which the dust cools predominantly. This in practice means that this should go all the way from
0.1 micron to 1000 micron
\end{sphinxadmonition}

\end{fulllineitems}

\index{get\_diskcontrol() (setup.radmc3d\_setup method)@\spxentry{get\_diskcontrol()}\spxextra{setup.radmc3d\_setup method}}

\begin{fulllineitems}
\phantomsection\label{\detokenize{radmc3d:setup.radmc3d_setup.get_diskcontrol}}
\pysigstartsignatures
\pysiglinewithargsret
{\sphinxbfcode{\sphinxupquote{get\_diskcontrol}}}
{\sphinxparam{\DUrole{n}{d\_to\_g\_ratio}\DUrole{o}{=}\DUrole{default_value}{0.01}}\sphinxparamcomma \sphinxparam{\DUrole{n}{a\_max}\DUrole{o}{=}\DUrole{default_value}{None}}\sphinxparamcomma \sphinxparam{\DUrole{n}{q}\DUrole{o}{=}\DUrole{default_value}{\sphinxhyphen{}3.5}}\sphinxparamcomma \sphinxparam{\DUrole{n}{Mass\_of\_star}\DUrole{o}{=}\DUrole{default_value}{None}}\sphinxparamcomma \sphinxparam{\DUrole{n}{Accretion\_rate}\DUrole{o}{=}\DUrole{default_value}{None}}\sphinxparamcomma \sphinxparam{\DUrole{n}{Radius\_of\_disk}\DUrole{o}{=}\DUrole{default_value}{None}}\sphinxparamcomma \sphinxparam{\DUrole{n}{Q}\DUrole{o}{=}\DUrole{default_value}{1.5}}\sphinxparamcomma \sphinxparam{\DUrole{n}{NR}\DUrole{o}{=}\DUrole{default_value}{None}}\sphinxparamcomma \sphinxparam{\DUrole{n}{NTheta}\DUrole{o}{=}\DUrole{default_value}{None}}\sphinxparamcomma \sphinxparam{\DUrole{n}{NPhi}\DUrole{o}{=}\DUrole{default_value}{None}}\sphinxparamcomma \sphinxparam{\DUrole{n}{generate\_table\_only}\DUrole{o}{=}\DUrole{default_value}{False}}}
{}
\pysigstopsignatures
\sphinxAtStartPar
Preparing the control file for disk model.
\begin{quote}\begin{description}
\sphinxlineitem{Parameters}\begin{itemize}
\item {} 
\sphinxAtStartPar
\sphinxstyleliteralstrong{\sphinxupquote{d\_to\_g\_ratio}} (\sphinxstyleliteralemphasis{\sphinxupquote{float}}) \textendash{} dust\sphinxhyphen{}to\sphinxhyphen{}gas mass ratio

\item {} 
\sphinxAtStartPar
\sphinxstyleliteralstrong{\sphinxupquote{a\_max}} (\sphinxstyleliteralemphasis{\sphinxupquote{float}}) \textendash{} Maxmum grain size (unit: mm)

\item {} 
\sphinxAtStartPar
\sphinxstyleliteralstrong{\sphinxupquote{q}} (\sphinxstyleliteralemphasis{\sphinxupquote{float}}) \textendash{} Slope for grain size distribution

\item {} 
\sphinxAtStartPar
\sphinxstyleliteralstrong{\sphinxupquote{Mass\_of\_star}} (\sphinxstyleliteralemphasis{\sphinxupquote{float}}) \textendash{} Mass of protostar (unit: M\_sun)

\item {} 
\sphinxAtStartPar
\sphinxstyleliteralstrong{\sphinxupquote{Accretion\_rate}} (\sphinxstyleliteralemphasis{\sphinxupquote{float}}) \textendash{} Accretion rate    (unit: M\_sun/yr)

\item {} 
\sphinxAtStartPar
\sphinxstyleliteralstrong{\sphinxupquote{Radius\_of\_disk}} (\sphinxstyleliteralemphasis{\sphinxupquote{float}}) \textendash{} Radius of disk    (unit: AU)

\item {} 
\sphinxAtStartPar
\sphinxstyleliteralstrong{\sphinxupquote{Q}} (\sphinxstyleliteralemphasis{\sphinxupquote{float}}) \textendash{} Toomre index

\item {} 
\sphinxAtStartPar
\sphinxstyleliteralstrong{\sphinxupquote{NR}} (\sphinxstyleliteralemphasis{\sphinxupquote{int}}) \textendash{} Resolution in R axis

\item {} 
\sphinxAtStartPar
\sphinxstyleliteralstrong{\sphinxupquote{NTheta}} (\sphinxstyleliteralemphasis{\sphinxupquote{int}}) \textendash{} Resolution in theta axis

\item {} 
\sphinxAtStartPar
\sphinxstyleliteralstrong{\sphinxupquote{NPhi}} (\sphinxstyleliteralemphasis{\sphinxupquote{int}}) \textendash{} Resolution in ohi axis

\end{itemize}

\end{description}\end{quote}

\end{fulllineitems}

\index{get\_gasdensitycontrol() (setup.radmc3d\_setup method)@\spxentry{get\_gasdensitycontrol()}\spxextra{setup.radmc3d\_setup method}}

\begin{fulllineitems}
\phantomsection\label{\detokenize{radmc3d:setup.radmc3d_setup.get_gasdensitycontrol}}
\pysigstartsignatures
\pysiglinewithargsret
{\sphinxbfcode{\sphinxupquote{get\_gasdensitycontrol}}}
{\sphinxparam{\DUrole{n}{abundance}\DUrole{o}{=}\DUrole{default_value}{1e\sphinxhyphen{}10}}\sphinxparamcomma \sphinxparam{\DUrole{n}{snowline}\DUrole{o}{=}\DUrole{default_value}{None}}\sphinxparamcomma \sphinxparam{\DUrole{n}{enhancement}\DUrole{o}{=}\DUrole{default_value}{100000.0}}\sphinxparamcomma \sphinxparam{\DUrole{n}{gas\_inside\_rcb}\DUrole{o}{=}\DUrole{default_value}{True}}\sphinxparamcomma \sphinxparam{\DUrole{o}{**}\DUrole{n}{kwargs}}}
{}
\pysigstopsignatures
\sphinxAtStartPar
Preparing the control file for gas density.
\begin{quote}\begin{description}
\sphinxlineitem{Parameters}\begin{itemize}
\item {} 
\sphinxAtStartPar
\sphinxstyleliteralstrong{\sphinxupquote{abundance}} (\sphinxstyleliteralemphasis{\sphinxupquote{float}}) \textendash{} Abundance of the simulated gas molecule compared with hydrogen

\item {} 
\sphinxAtStartPar
\sphinxstyleliteralstrong{\sphinxupquote{snowline}} (\sphinxstyleliteralemphasis{\sphinxupquote{float}}) \textendash{} At which temperature causes desorption from dust grains (unit:K)
(if None, there is no abundance enhancement)

\item {} 
\sphinxAtStartPar
\sphinxstyleliteralstrong{\sphinxupquote{enhancement}} (\sphinxstyleliteralemphasis{\sphinxupquote{float}}) \textendash{} How much the abundance is enhanced inside snowline

\item {} 
\sphinxAtStartPar
\sphinxstyleliteralstrong{\sphinxupquote{gas\_inside\_rcb}} (\sphinxstyleliteralemphasis{\sphinxupquote{bool}}) \textendash{} Whether gas is absent inside centrifugal barrier

\end{itemize}

\end{description}\end{quote}

\end{fulllineitems}

\index{get\_heatcontrol() (setup.radmc3d\_setup method)@\spxentry{get\_heatcontrol()}\spxextra{setup.radmc3d\_setup method}}

\begin{fulllineitems}
\phantomsection\label{\detokenize{radmc3d:setup.radmc3d_setup.get_heatcontrol}}
\pysigstartsignatures
\pysiglinewithargsret
{\sphinxbfcode{\sphinxupquote{get\_heatcontrol}}}
{\sphinxparam{\DUrole{n}{L\_star}\DUrole{o}{=}\DUrole{default_value}{None}}\sphinxparamcomma \sphinxparam{\DUrole{n}{R\_star}\DUrole{o}{=}\DUrole{default_value}{None}}\sphinxparamcomma \sphinxparam{\DUrole{n}{accretion}\DUrole{o}{=}\DUrole{default_value}{True}}\sphinxparamcomma \sphinxparam{\DUrole{n}{irradiation}\DUrole{o}{=}\DUrole{default_value}{True}}\sphinxparamcomma \sphinxparam{\DUrole{o}{**}\DUrole{n}{kwargs}}}
{}
\pysigstopsignatures
\sphinxAtStartPar
Preparing the control file for heating mechanism.
\begin{quote}\begin{description}
\sphinxlineitem{Parameters}\begin{itemize}
\item {} 
\sphinxAtStartPar
\sphinxstyleliteralstrong{\sphinxupquote{accretion}} (\sphinxstyleliteralemphasis{\sphinxupquote{bool}}) \textendash{} Accretion heating mechanism due to release of gravitational energy

\item {} 
\sphinxAtStartPar
\sphinxstyleliteralstrong{\sphinxupquote{irradiation}} (\sphinxstyleliteralemphasis{\sphinxupquote{bool}}) \textendash{} Irradiation heating from central protostar

\item {} 
\sphinxAtStartPar
\sphinxstyleliteralstrong{\sphinxupquote{L\_star}} (\sphinxstyleliteralemphasis{\sphinxupquote{float}}) \textendash{} Luminosity of central protostar (unit: L\_sun)

\end{itemize}

\end{description}\end{quote}

\begin{sphinxadmonition}{note}{Note:}
\sphinxAtStartPar
If both heating mechanisms are True, combine accretion and irradiation temperature map.
\end{sphinxadmonition}

\end{fulllineitems}

\index{get\_linecontrol() (setup.radmc3d\_setup method)@\spxentry{get\_linecontrol()}\spxextra{setup.radmc3d\_setup method}}

\begin{fulllineitems}
\phantomsection\label{\detokenize{radmc3d:setup.radmc3d_setup.get_linecontrol}}
\pysigstartsignatures
\pysiglinewithargsret
{\sphinxbfcode{\sphinxupquote{get\_linecontrol}}}
{\sphinxparam{\DUrole{n}{filename}\DUrole{o}{=}\DUrole{default_value}{None}}\sphinxparamcomma \sphinxparam{\DUrole{n}{comment}\DUrole{o}{=}\DUrole{default_value}{None}}\sphinxparamcomma \sphinxparam{\DUrole{o}{**}\DUrole{n}{kwargs}}}
{}
\pysigstopsignatures
\sphinxAtStartPar
Preparing the control file for lines.

\sphinxAtStartPar
Input :
\begin{description}
\sphinxlineitem{filename (string)}{[}output filename. (default: lines.inp).{]}
\sphinxAtStartPar
It will still creat a file with default name,
but will duplicate an output file with the specified name.

\end{description}

\sphinxAtStartPar
See example for how to include lines

\end{fulllineitems}

\index{get\_mastercontrol() (setup.radmc3d\_setup method)@\spxentry{get\_mastercontrol()}\spxextra{setup.radmc3d\_setup method}}

\begin{fulllineitems}
\phantomsection\label{\detokenize{radmc3d:setup.radmc3d_setup.get_mastercontrol}}
\pysigstartsignatures
\pysiglinewithargsret
{\sphinxbfcode{\sphinxupquote{get\_mastercontrol}}}
{\sphinxparam{\DUrole{n}{filename}\DUrole{o}{=}\DUrole{default_value}{None}}\sphinxparamcomma \sphinxparam{\DUrole{n}{comment}\DUrole{o}{=}\DUrole{default_value}{None}}\sphinxparamcomma \sphinxparam{\DUrole{n}{incl\_dust}\DUrole{o}{=}\DUrole{default_value}{None}}\sphinxparamcomma \sphinxparam{\DUrole{n}{incl\_lines}\DUrole{o}{=}\DUrole{default_value}{None}}\sphinxparamcomma \sphinxparam{\DUrole{n}{nphot}\DUrole{o}{=}\DUrole{default_value}{1000000}}\sphinxparamcomma \sphinxparam{\DUrole{n}{nphot\_scat}\DUrole{o}{=}\DUrole{default_value}{1000000}}\sphinxparamcomma \sphinxparam{\DUrole{n}{scattering\_mode\_max}\DUrole{o}{=}\DUrole{default_value}{None}}\sphinxparamcomma \sphinxparam{\DUrole{n}{istar\_sphere}\DUrole{o}{=}\DUrole{default_value}{1}}\sphinxparamcomma \sphinxparam{\DUrole{n}{num\_cpu}\DUrole{o}{=}\DUrole{default_value}{None}}\sphinxparamcomma \sphinxparam{\DUrole{o}{**}\DUrole{n}{kwargs}}}
{}
\pysigstopsignatures
\sphinxAtStartPar
Preparing the master control file for radmc3d.
\subsubsection*{Example}

\sphinxAtStartPar
test.get\_mastercontrol(comment = ‘this is a test’, a=1.0, b=2.0, c=3.0)
\begin{quote}\begin{description}
\sphinxlineitem{Parameters}\begin{itemize}
\item {} 
\sphinxAtStartPar
\sphinxstyleliteralstrong{\sphinxupquote{filename}} (\sphinxstyleliteralemphasis{\sphinxupquote{string}}) \textendash{} output filename. (default: radmc3d.inp)
It will still creat a file with default name, but will duplicate an output file with the specified name.

\item {} 
\sphinxAtStartPar
\sphinxstyleliteralstrong{\sphinxupquote{comment}} (\sphinxstyleliteralemphasis{\sphinxupquote{string}}) \textendash{} comment to add to the file header (default: None)

\item {} 
\sphinxAtStartPar
\sphinxstyleliteralstrong{\sphinxupquote{incl\_dust}} (\sphinxstyleliteralemphasis{\sphinxupquote{0/1/None}}) \textendash{} 0: force not include dust/ 1: force include / None: let radmc3d determine

\item {} 
\sphinxAtStartPar
\sphinxstyleliteralstrong{\sphinxupquote{incl\_lines}} (\sphinxstyleliteralemphasis{\sphinxupquote{0/1/None}}) \textendash{} 0: force not include line/ 1: force include / None: let radmc3d determine

\item {} 
\sphinxAtStartPar
\sphinxstyleliteralstrong{\sphinxupquote{nphot}} (\sphinxstyleliteralemphasis{\sphinxupquote{int}}) \textendash{} The number of photon packages used for the thermal Monte Carlo simulation (default: 1000000)

\item {} 
\sphinxAtStartPar
\sphinxstyleliteralstrong{\sphinxupquote{nphot\_scat}} (\sphinxstyleliteralemphasis{\sphinxupquote{int}}) \textendash{} The number of photon packages for the scattering Monte Carlo simulations, done before image\sphinxhyphen{}rendering (default: 1000000)

\item {} 
\sphinxAtStartPar
\sphinxstyleliteralstrong{\sphinxupquote{scattering\_mode\_max}} (\sphinxstyleliteralemphasis{\sphinxupquote{0/1/2/None}}) \textendash{} 0: no scattering / 1: isotropic scattering / 2: full scattering / None: let radmc decide (default: None)

\item {} 
\sphinxAtStartPar
\sphinxstyleliteralstrong{\sphinxupquote{istar\_sphere}} (\sphinxstyleliteralemphasis{\sphinxupquote{0/1}}) \textendash{} If 0/1, treat stars as point\sphinxhyphen{}source/sphere (default: 1)

\item {} 
\sphinxAtStartPar
\sphinxstyleliteralstrong{\sphinxupquote{num\_cpu}} (\sphinxstyleliteralemphasis{\sphinxupquote{int}}) \textendash{} number of cpu core to use (default: available threads\sphinxhyphen{}2)

\end{itemize}

\end{description}\end{quote}

\begin{sphinxadmonition}{note}{Note:}\begin{description}
\sphinxlineitem{Other possible options (including using {\color{red}\bfseries{}**}kwargs) see}
\sphinxAtStartPar
\sphinxurl{https://www.ita.uni-heidelberg.de/~dullemond/software/radmc-3d/manual\_radmc3d/inputoutputfiles.html}

\end{description}
\end{sphinxadmonition}

\end{fulllineitems}

\index{get\_vfieldcontrol() (setup.radmc3d\_setup method)@\spxentry{get\_vfieldcontrol()}\spxextra{setup.radmc3d\_setup method}}

\begin{fulllineitems}
\phantomsection\label{\detokenize{radmc3d:setup.radmc3d_setup.get_vfieldcontrol}}
\pysigstartsignatures
\pysiglinewithargsret
{\sphinxbfcode{\sphinxupquote{get\_vfieldcontrol}}}
{\sphinxparam{\DUrole{n}{Kep}\DUrole{o}{=}\DUrole{default_value}{True}}\sphinxparamcomma \sphinxparam{\DUrole{n}{vinfall}\DUrole{o}{=}\DUrole{default_value}{None}}\sphinxparamcomma \sphinxparam{\DUrole{n}{Rcb}\DUrole{o}{=}\DUrole{default_value}{None}}\sphinxparamcomma \sphinxparam{\DUrole{n}{outflow}\DUrole{o}{=}\DUrole{default_value}{None}}}
{}
\pysigstopsignatures
\sphinxAtStartPar
Preparing the control file for velocity field.
\begin{quote}\begin{description}
\sphinxlineitem{Parameters}\begin{itemize}
\item {} 
\sphinxAtStartPar
\sphinxstyleliteralstrong{\sphinxupquote{Kep}} (\sphinxstyleliteralemphasis{\sphinxupquote{bool}}) \textendash{} Keplerian azimuthal velocity field

\item {} 
\sphinxAtStartPar
\sphinxstyleliteralstrong{\sphinxupquote{vinfall}} (\sphinxstyleliteralemphasis{\sphinxupquote{float}}) \textendash{} Infall velocity (unit: Keplerian velocity) e.g., 1 for 1 Keplerian velocity of infall direction

\item {} 
\sphinxAtStartPar
\sphinxstyleliteralstrong{\sphinxupquote{Rcb}} (\sphinxstyleliteralemphasis{\sphinxupquote{float}}) \textendash{} Centrifugal barrier (unit: AU)

\item {} 
\sphinxAtStartPar
\sphinxstyleliteralstrong{\sphinxupquote{outflow}} (\sphinxstyleliteralemphasis{\sphinxupquote{float}}) \textendash{} Outflow velocity (unit: Keplerian velocity)

\end{itemize}

\end{description}\end{quote}

\end{fulllineitems}

\index{write\_amr\_grid() (setup.radmc3d\_setup method)@\spxentry{write\_amr\_grid()}\spxextra{setup.radmc3d\_setup method}}

\begin{fulllineitems}
\phantomsection\label{\detokenize{radmc3d:setup.radmc3d_setup.write_amr_grid}}
\pysigstartsignatures
\pysiglinewithargsret
{\sphinxbfcode{\sphinxupquote{write\_amr\_grid}}}
{}
{}
\pysigstopsignatures
\end{fulllineitems}

\index{write\_dust\_density() (setup.radmc3d\_setup method)@\spxentry{write\_dust\_density()}\spxextra{setup.radmc3d\_setup method}}

\begin{fulllineitems}
\phantomsection\label{\detokenize{radmc3d:setup.radmc3d_setup.write_dust_density}}
\pysigstartsignatures
\pysiglinewithargsret
{\sphinxbfcode{\sphinxupquote{write\_dust\_density}}}
{}
{}
\pysigstopsignatures
\sphinxAtStartPar
Preparing the control file for dust density.

\end{fulllineitems}

\index{write\_dust\_opac() (setup.radmc3d\_setup method)@\spxentry{write\_dust\_opac()}\spxextra{setup.radmc3d\_setup method}}

\begin{fulllineitems}
\phantomsection\label{\detokenize{radmc3d:setup.radmc3d_setup.write_dust_opac}}
\pysigstartsignatures
\pysiglinewithargsret
{\sphinxbfcode{\sphinxupquote{write\_dust\_opac}}}
{}
{}
\pysigstopsignatures
\sphinxAtStartPar
Preparing the control file for dust opacity.

\end{fulllineitems}


\end{fulllineitems}

\index{generate\_simulation (class in simulate)@\spxentry{generate\_simulation}\spxextra{class in simulate}}

\begin{fulllineitems}
\phantomsection\label{\detokenize{radmc3d:simulate.generate_simulation}}
\pysigstartsignatures
\pysiglinewithargsret
{\sphinxbfcode{\sphinxupquote{\DUrole{k}{class}\DUrole{w}{ }}}\sphinxcode{\sphinxupquote{simulate.}}\sphinxbfcode{\sphinxupquote{generate\_simulation}}}
{\sphinxparam{\DUrole{n}{save\_npz}\DUrole{o}{=}\DUrole{default_value}{True}}\sphinxparamcomma \sphinxparam{\DUrole{n}{save\_out}\DUrole{o}{=}\DUrole{default_value}{True}}}
{}
\pysigstopsignatures
\sphinxAtStartPar
Bases: \sphinxcode{\sphinxupquote{object}}
\subsubsection*{Methods}


\begin{savenotes}
\sphinxatlongtablestart
\sphinxthistablewithglobalstyle
\sphinxthistablewithnovlinesstyle
\makeatletter
  \LTleft \@totalleftmargin plus1fill
  \LTright\dimexpr\columnwidth-\@totalleftmargin-\linewidth\relax plus1fill
\makeatother
\begin{longtable}{\X{1}{2}\X{1}{2}}
\sphinxtoprule
\endfirsthead

\multicolumn{2}{c}{\sphinxnorowcolor
    \makebox[0pt]{\sphinxtablecontinued{\tablename\ \thetable{} \textendash{} continued from previous page}}%
}\\
\sphinxtoprule
\endhead

\sphinxbottomrule
\multicolumn{2}{r}{\sphinxnorowcolor
    \makebox[0pt][r]{\sphinxtablecontinued{continues on next page}}%
}\\
\endfoot

\endlastfoot
\sphinxtableatstartofbodyhook

\sphinxAtStartPar
{\hyperref[\detokenize{radmc3d:simulate.generate_simulation.generate_continuum}]{\sphinxcrossref{\sphinxcode{\sphinxupquote{generate\_continuum}}}}}({[}dir, fname, incl, wav, ...{]})
&
\sphinxAtStartPar
This function will generate a continuum image.
\\
\sphinxhline
\sphinxAtStartPar
{\hyperref[\detokenize{radmc3d:simulate.generate_simulation.generate_cube}]{\sphinxcrossref{\sphinxcode{\sphinxupquote{generate\_cube}}}}}({[}dir, fname, npix, sizeau, ...{]})
&
\sphinxAtStartPar
This function will generate a cube of spectral line.
\\
\sphinxhline
\sphinxAtStartPar
{\hyperref[\detokenize{radmc3d:simulate.generate_simulation.generate_line_spectrum}]{\sphinxcrossref{\sphinxcode{\sphinxupquote{generate\_line\_spectrum}}}}}({[}dir, fname, incl, ...{]})
&
\sphinxAtStartPar
This function will generate a line spectrum.
\\
\sphinxhline
\sphinxAtStartPar
{\hyperref[\detokenize{radmc3d:simulate.generate_simulation.generate_sed}]{\sphinxcrossref{\sphinxcode{\sphinxupquote{generate\_sed}}}}}({[}dir, fname, incl, freq\_min, ...{]})
&
\sphinxAtStartPar
This function will generate a spectral energy distribution (SED).
\\
\sphinxhline
\sphinxAtStartPar
{\hyperref[\detokenize{radmc3d:simulate.generate_simulation.save_npzfile}]{\sphinxcrossref{\sphinxcode{\sphinxupquote{save\_npzfile}}}}}(data, dir, fname, **kwargs)
&
\sphinxAtStartPar
This will only save image data regardless of other information The file tyep will be \textquotesingle{}{\color{red}\bfseries{}*}.npz\textquotesingle{}, which the storage may be kB\sphinxhyphen{}scale
\\
\sphinxhline
\sphinxAtStartPar
{\hyperref[\detokenize{radmc3d:simulate.generate_simulation.save_outfile}]{\sphinxcrossref{\sphinxcode{\sphinxupquote{save\_outfile}}}}}(dir, fname, **kwargs)
&
\sphinxAtStartPar
This will save whole information of simulation. The file type will be \textquotesingle{}{\color{red}\bfseries{}*}.out\textquotesingle{}, which the storage may be MB\sphinxhyphen{}scale.
\\
\sphinxbottomrule
\end{longtable}
\sphinxtableafterendhook
\sphinxatlongtableend
\end{savenotes}
\index{generate\_continuum() (simulate.generate\_simulation method)@\spxentry{generate\_continuum()}\spxextra{simulate.generate\_simulation method}}

\begin{fulllineitems}
\phantomsection\label{\detokenize{radmc3d:simulate.generate_simulation.generate_continuum}}
\pysigstartsignatures
\pysiglinewithargsret
{\sphinxbfcode{\sphinxupquote{generate\_continuum}}}
{\sphinxparam{\DUrole{n}{dir}\DUrole{o}{=}\DUrole{default_value}{\textquotesingle{}./test/\textquotesingle{}}}\sphinxparamcomma \sphinxparam{\DUrole{n}{fname}\DUrole{o}{=}\DUrole{default_value}{\textquotesingle{}test\textquotesingle{}}}\sphinxparamcomma \sphinxparam{\DUrole{n}{incl}\DUrole{o}{=}\DUrole{default_value}{73}}\sphinxparamcomma \sphinxparam{\DUrole{n}{wav}\DUrole{o}{=}\DUrole{default_value}{1300}}\sphinxparamcomma \sphinxparam{\DUrole{n}{npix}\DUrole{o}{=}\DUrole{default_value}{200}}\sphinxparamcomma \sphinxparam{\DUrole{n}{sizeau}\DUrole{o}{=}\DUrole{default_value}{100}}\sphinxparamcomma \sphinxparam{\DUrole{n}{posang}\DUrole{o}{=}\DUrole{default_value}{45}}\sphinxparamcomma \sphinxparam{\DUrole{n}{scat}\DUrole{o}{=}\DUrole{default_value}{True}}\sphinxparamcomma \sphinxparam{\DUrole{o}{**}\DUrole{n}{kwargs}}}
{}
\pysigstopsignatures
\sphinxAtStartPar
This function will generate a continuum image.
\begin{quote}\begin{description}
\sphinxlineitem{Parameters}\begin{itemize}
\item {} 
\sphinxAtStartPar
\sphinxstyleliteralstrong{\sphinxupquote{dir}} (\sphinxstyleliteralemphasis{\sphinxupquote{str}}) \textendash{} Directory to save the output

\item {} 
\sphinxAtStartPar
\sphinxstyleliteralstrong{\sphinxupquote{fname}} (\sphinxstyleliteralemphasis{\sphinxupquote{str}}) \textendash{} File name to save the output

\item {} 
\sphinxAtStartPar
\sphinxstyleliteralstrong{\sphinxupquote{incl}} (\sphinxstyleliteralemphasis{\sphinxupquote{float}}) \textendash{} Inclination angle of the disk

\item {} 
\sphinxAtStartPar
\sphinxstyleliteralstrong{\sphinxupquote{wav}} (\sphinxstyleliteralemphasis{\sphinxupquote{float}}) \textendash{} Wavelength to simulate

\item {} 
\sphinxAtStartPar
\sphinxstyleliteralstrong{\sphinxupquote{npix}} (\sphinxstyleliteralemphasis{\sphinxupquote{int}}) \textendash{} Number of map’s pixels

\item {} 
\sphinxAtStartPar
\sphinxstyleliteralstrong{\sphinxupquote{sizeau}} (\sphinxstyleliteralemphasis{\sphinxupquote{float}}) \textendash{} Map’s span

\item {} 
\sphinxAtStartPar
\sphinxstyleliteralstrong{\sphinxupquote{posang}} (\sphinxstyleliteralemphasis{\sphinxupquote{float}}) \textendash{} Position angle of the disk

\item {} 
\sphinxAtStartPar
\sphinxstyleliteralstrong{\sphinxupquote{scat}} (\sphinxstyleliteralemphasis{\sphinxupquote{bool}}) \textendash{} If True, scattering is included. (Time\sphinxhyphen{}consuming)

\end{itemize}

\end{description}\end{quote}

\end{fulllineitems}

\index{generate\_cube() (simulate.generate\_simulation method)@\spxentry{generate\_cube()}\spxextra{simulate.generate\_simulation method}}

\begin{fulllineitems}
\phantomsection\label{\detokenize{radmc3d:simulate.generate_simulation.generate_cube}}
\pysigstartsignatures
\pysiglinewithargsret
{\sphinxbfcode{\sphinxupquote{generate\_cube}}}
{\sphinxparam{\DUrole{n}{dir}\DUrole{o}{=}\DUrole{default_value}{\textquotesingle{}./test/\textquotesingle{}}}\sphinxparamcomma \sphinxparam{\DUrole{n}{fname}\DUrole{o}{=}\DUrole{default_value}{\textquotesingle{}test\textquotesingle{}}}\sphinxparamcomma \sphinxparam{\DUrole{n}{npix}\DUrole{o}{=}\DUrole{default_value}{100}}\sphinxparamcomma \sphinxparam{\DUrole{n}{sizeau}\DUrole{o}{=}\DUrole{default_value}{100}}\sphinxparamcomma \sphinxparam{\DUrole{n}{incl}\DUrole{o}{=}\DUrole{default_value}{73}}\sphinxparamcomma \sphinxparam{\DUrole{n}{line}\DUrole{o}{=}\DUrole{default_value}{240}}\sphinxparamcomma \sphinxparam{\DUrole{n}{v\_width}\DUrole{o}{=}\DUrole{default_value}{10}}\sphinxparamcomma \sphinxparam{\DUrole{n}{vkms}\DUrole{o}{=}\DUrole{default_value}{0}}\sphinxparamcomma \sphinxparam{\DUrole{n}{nlam}\DUrole{o}{=}\DUrole{default_value}{11}}\sphinxparamcomma \sphinxparam{\DUrole{n}{posang}\DUrole{o}{=}\DUrole{default_value}{45}}\sphinxparamcomma \sphinxparam{\DUrole{n}{nodust}\DUrole{o}{=}\DUrole{default_value}{False}}\sphinxparamcomma \sphinxparam{\DUrole{n}{scat}\DUrole{o}{=}\DUrole{default_value}{True}}\sphinxparamcomma \sphinxparam{\DUrole{n}{extract\_gas}\DUrole{o}{=}\DUrole{default_value}{True}}\sphinxparamcomma \sphinxparam{\DUrole{o}{**}\DUrole{n}{kwargs}}}
{}
\pysigstopsignatures
\sphinxAtStartPar
This function will generate a cube of spectral line.
\begin{quote}\begin{description}
\sphinxlineitem{Parameters}\begin{itemize}
\item {} 
\sphinxAtStartPar
\sphinxstyleliteralstrong{\sphinxupquote{dir}} (\sphinxstyleliteralemphasis{\sphinxupquote{str}}) \textendash{} Directory to save the output

\item {} 
\sphinxAtStartPar
\sphinxstyleliteralstrong{\sphinxupquote{fname}} (\sphinxstyleliteralemphasis{\sphinxupquote{str}}) \textendash{} File name to save the output

\item {} 
\sphinxAtStartPar
\sphinxstyleliteralstrong{\sphinxupquote{npix}} (\sphinxstyleliteralemphasis{\sphinxupquote{int}}) \textendash{} Number of map’s pixels

\item {} 
\sphinxAtStartPar
\sphinxstyleliteralstrong{\sphinxupquote{sizeau}} (\sphinxstyleliteralemphasis{\sphinxupquote{float}}) \textendash{} Map’s span

\item {} 
\sphinxAtStartPar
\sphinxstyleliteralstrong{\sphinxupquote{incl}} (\sphinxstyleliteralemphasis{\sphinxupquote{float}}) \textendash{} Inclination angle of the disk

\item {} 
\sphinxAtStartPar
\sphinxstyleliteralstrong{\sphinxupquote{line}} (\sphinxstyleliteralemphasis{\sphinxupquote{int}}) \textendash{} Transistion level (see ‘molecule\_ch3oh.inp’)

\item {} 
\sphinxAtStartPar
\sphinxstyleliteralstrong{\sphinxupquote{v\_width}} (\sphinxstyleliteralemphasis{\sphinxupquote{float}}) \textendash{} Range of velocity to simulate

\item {} 
\sphinxAtStartPar
\sphinxstyleliteralstrong{\sphinxupquote{vkms}} (\sphinxstyleliteralemphasis{\sphinxupquote{float}}) \textendash{} Velocity of the line center

\item {} 
\sphinxAtStartPar
\sphinxstyleliteralstrong{\sphinxupquote{nlam}} (\sphinxstyleliteralemphasis{\sphinxupquote{int}}) \textendash{} Number of velocities

\item {} 
\sphinxAtStartPar
\sphinxstyleliteralstrong{\sphinxupquote{nodust}} (\sphinxstyleliteralemphasis{\sphinxupquote{bool}}) \textendash{} If False, dust effect is included

\item {} 
\sphinxAtStartPar
\sphinxstyleliteralstrong{\sphinxupquote{scat}} (\sphinxstyleliteralemphasis{\sphinxupquote{bool}}) \textendash{} If True and nodust=False, scattering is included. (Time\sphinxhyphen{}consuming)

\item {} 
\sphinxAtStartPar
\sphinxstyleliteralstrong{\sphinxupquote{extracted\_gas}} (\sphinxstyleliteralemphasis{\sphinxupquote{bool}}) \textendash{} If True, spectral line is extracted (I\_\{dust+gas\}\sphinxhyphen{}I\_\{dust\})

\end{itemize}

\end{description}\end{quote}

\end{fulllineitems}

\index{generate\_line\_spectrum() (simulate.generate\_simulation method)@\spxentry{generate\_line\_spectrum()}\spxextra{simulate.generate\_simulation method}}

\begin{fulllineitems}
\phantomsection\label{\detokenize{radmc3d:simulate.generate_simulation.generate_line_spectrum}}
\pysigstartsignatures
\pysiglinewithargsret
{\sphinxbfcode{\sphinxupquote{generate\_line\_spectrum}}}
{\sphinxparam{\DUrole{n}{dir}\DUrole{o}{=}\DUrole{default_value}{\textquotesingle{}./test/\textquotesingle{}}}\sphinxparamcomma \sphinxparam{\DUrole{n}{fname}\DUrole{o}{=}\DUrole{default_value}{\textquotesingle{}test\textquotesingle{}}}\sphinxparamcomma \sphinxparam{\DUrole{n}{incl}\DUrole{o}{=}\DUrole{default_value}{73}}\sphinxparamcomma \sphinxparam{\DUrole{n}{line}\DUrole{o}{=}\DUrole{default_value}{240}}\sphinxparamcomma \sphinxparam{\DUrole{n}{v\_width}\DUrole{o}{=}\DUrole{default_value}{10}}\sphinxparamcomma \sphinxparam{\DUrole{n}{nlam}\DUrole{o}{=}\DUrole{default_value}{10}}\sphinxparamcomma \sphinxparam{\DUrole{n}{vkms}\DUrole{o}{=}\DUrole{default_value}{0}}\sphinxparamcomma \sphinxparam{\DUrole{n}{nodust}\DUrole{o}{=}\DUrole{default_value}{False}}\sphinxparamcomma \sphinxparam{\DUrole{n}{scat}\DUrole{o}{=}\DUrole{default_value}{True}}\sphinxparamcomma \sphinxparam{\DUrole{n}{extract\_gas}\DUrole{o}{=}\DUrole{default_value}{True}}\sphinxparamcomma \sphinxparam{\DUrole{o}{**}\DUrole{n}{kwargs}}}
{}
\pysigstopsignatures
\sphinxAtStartPar
This function will generate a line spectrum.
\begin{quote}\begin{description}
\sphinxlineitem{Parameters}\begin{itemize}
\item {} 
\sphinxAtStartPar
\sphinxstyleliteralstrong{\sphinxupquote{dir}} (\sphinxstyleliteralemphasis{\sphinxupquote{str}}) \textendash{} Directory to save the output

\item {} 
\sphinxAtStartPar
\sphinxstyleliteralstrong{\sphinxupquote{fname}} (\sphinxstyleliteralemphasis{\sphinxupquote{str}}) \textendash{} File name to save the output

\item {} 
\sphinxAtStartPar
\sphinxstyleliteralstrong{\sphinxupquote{incl}} (\sphinxstyleliteralemphasis{\sphinxupquote{float}}) \textendash{} Inclination angle of the disk

\item {} 
\sphinxAtStartPar
\sphinxstyleliteralstrong{\sphinxupquote{line}} (\sphinxstyleliteralemphasis{\sphinxupquote{int}}) \textendash{} Transistion level (see ‘molecule\_ch3oh.inp’)

\item {} 
\sphinxAtStartPar
\sphinxstyleliteralstrong{\sphinxupquote{v\_width}} (\sphinxstyleliteralemphasis{\sphinxupquote{float}}) \textendash{} Range of velocity to simulate

\item {} 
\sphinxAtStartPar
\sphinxstyleliteralstrong{\sphinxupquote{nlam}} (\sphinxstyleliteralemphasis{\sphinxupquote{int}}) \textendash{} Number of velocities

\item {} 
\sphinxAtStartPar
\sphinxstyleliteralstrong{\sphinxupquote{vkms}} (\sphinxstyleliteralemphasis{\sphinxupquote{float}}) \textendash{} Velocity of the line center

\item {} 
\sphinxAtStartPar
\sphinxstyleliteralstrong{\sphinxupquote{nodust}} (\sphinxstyleliteralemphasis{\sphinxupquote{bool}}) \textendash{} If False, dust effect is included

\item {} 
\sphinxAtStartPar
\sphinxstyleliteralstrong{\sphinxupquote{scat}} (\sphinxstyleliteralemphasis{\sphinxupquote{bool}}) \textendash{} If True and nodust=False, scattering is included. (Time\sphinxhyphen{}consuming)

\item {} 
\sphinxAtStartPar
\sphinxstyleliteralstrong{\sphinxupquote{extracted\_gas}} (\sphinxstyleliteralemphasis{\sphinxupquote{bool}}) \textendash{} If True, spectral line is extracted (I\_\{dust+gas\}\sphinxhyphen{}I\_\{dust\})

\end{itemize}

\end{description}\end{quote}

\end{fulllineitems}

\index{generate\_sed() (simulate.generate\_simulation method)@\spxentry{generate\_sed()}\spxextra{simulate.generate\_simulation method}}

\begin{fulllineitems}
\phantomsection\label{\detokenize{radmc3d:simulate.generate_simulation.generate_sed}}
\pysigstartsignatures
\pysiglinewithargsret
{\sphinxbfcode{\sphinxupquote{generate\_sed}}}
{\sphinxparam{\DUrole{n}{dir}\DUrole{o}{=}\DUrole{default_value}{\textquotesingle{}./test/\textquotesingle{}}}\sphinxparamcomma \sphinxparam{\DUrole{n}{fname}\DUrole{o}{=}\DUrole{default_value}{\textquotesingle{}test\textquotesingle{}}}\sphinxparamcomma \sphinxparam{\DUrole{n}{incl}\DUrole{o}{=}\DUrole{default_value}{73}}\sphinxparamcomma \sphinxparam{\DUrole{n}{freq\_min}\DUrole{o}{=}\DUrole{default_value}{50.0}}\sphinxparamcomma \sphinxparam{\DUrole{n}{freq\_max}\DUrole{o}{=}\DUrole{default_value}{1000.0}}\sphinxparamcomma \sphinxparam{\DUrole{n}{nlam}\DUrole{o}{=}\DUrole{default_value}{100}}\sphinxparamcomma \sphinxparam{\DUrole{n}{scat}\DUrole{o}{=}\DUrole{default_value}{True}}\sphinxparamcomma \sphinxparam{\DUrole{o}{**}\DUrole{n}{kwargs}}}
{}
\pysigstopsignatures
\sphinxAtStartPar
This function will generate a spectral energy distribution (SED).
\begin{quote}\begin{description}
\sphinxlineitem{Parameters}\begin{itemize}
\item {} 
\sphinxAtStartPar
\sphinxstyleliteralstrong{\sphinxupquote{dir}} (\sphinxstyleliteralemphasis{\sphinxupquote{str}}) \textendash{} Directory to save the output

\item {} 
\sphinxAtStartPar
\sphinxstyleliteralstrong{\sphinxupquote{fname}} (\sphinxstyleliteralemphasis{\sphinxupquote{str}}) \textendash{} File name to save the output

\item {} 
\sphinxAtStartPar
\sphinxstyleliteralstrong{\sphinxupquote{incl}} (\sphinxstyleliteralemphasis{\sphinxupquote{float}}) \textendash{} Inclination angle of the disk

\item {} 
\sphinxAtStartPar
\sphinxstyleliteralstrong{\sphinxupquote{freq\_min}} (\sphinxstyleliteralemphasis{\sphinxupquote{float}}) \textendash{} Minimum frequency to simulate

\item {} 
\sphinxAtStartPar
\sphinxstyleliteralstrong{\sphinxupquote{freq\_max}} (\sphinxstyleliteralemphasis{\sphinxupquote{float}}) \textendash{} Maximum frequency to simulate

\item {} 
\sphinxAtStartPar
\sphinxstyleliteralstrong{\sphinxupquote{nlam}} (\sphinxstyleliteralemphasis{\sphinxupquote{int}}) \textendash{} Number of wavelengths

\item {} 
\sphinxAtStartPar
\sphinxstyleliteralstrong{\sphinxupquote{scat}} (\sphinxstyleliteralemphasis{\sphinxupquote{bool}}) \textendash{} If True, scattering is included. (Time\sphinxhyphen{}consuming)

\end{itemize}

\end{description}\end{quote}

\end{fulllineitems}

\index{save\_npzfile() (simulate.generate\_simulation method)@\spxentry{save\_npzfile()}\spxextra{simulate.generate\_simulation method}}

\begin{fulllineitems}
\phantomsection\label{\detokenize{radmc3d:simulate.generate_simulation.save_npzfile}}
\pysigstartsignatures
\pysiglinewithargsret
{\sphinxbfcode{\sphinxupquote{save\_npzfile}}}
{\sphinxparam{\DUrole{n}{data}}\sphinxparamcomma \sphinxparam{\DUrole{n}{dir}}\sphinxparamcomma \sphinxparam{\DUrole{n}{fname}}\sphinxparamcomma \sphinxparam{\DUrole{o}{**}\DUrole{n}{kwargs}}}
{}
\pysigstopsignatures
\sphinxAtStartPar
This will only save image data regardless of other information
The file tyep will be ‘{\color{red}\bfseries{}*}.npz’, which the storage may be kB\sphinxhyphen{}scale

\end{fulllineitems}

\index{save\_outfile() (simulate.generate\_simulation method)@\spxentry{save\_outfile()}\spxextra{simulate.generate\_simulation method}}

\begin{fulllineitems}
\phantomsection\label{\detokenize{radmc3d:simulate.generate_simulation.save_outfile}}
\pysigstartsignatures
\pysiglinewithargsret
{\sphinxbfcode{\sphinxupquote{save\_outfile}}}
{\sphinxparam{\DUrole{n}{dir}}\sphinxparamcomma \sphinxparam{\DUrole{n}{fname}}\sphinxparamcomma \sphinxparam{\DUrole{o}{**}\DUrole{n}{kwargs}}}
{}
\pysigstopsignatures
\sphinxAtStartPar
This will save whole information of simulation.
The file type will be ‘{\color{red}\bfseries{}*}.out’, which the storage may be MB\sphinxhyphen{}scale

\end{fulllineitems}


\end{fulllineitems}


\sphinxstepscope


\chapter{Examples}
\label{\detokenize{example:examples}}\label{\detokenize{example::doc}}


\renewcommand{\indexname}{Index}
\printindex
\end{document}